% !TeX program = xelatex
% 编译：latexmk -xelatex -bibtex main.tex
% 官方硬约束已固化：摘要独占第1页；正文第2页开始；不生成目录；页码居中；
% AI工具使用声明位于参考文献之前；附录列支撑材料与源代码。
\documentclass[UTF8,a4paper,zihao=5]{ctexart}

\usepackage[left=2.5cm,right=2.5cm,top=2.5cm,bottom=2.3cm,
            headheight=14pt,footskip=1.2cm]{geometry}
\usepackage{amsmath,amssymb,mathtools}
\usepackage{booktabs,longtable,tabularx,array,multirow,threeparttable}
\usepackage{graphicx}
\usepackage{caption,subcaption}
\usepackage{xcolor}
\usepackage{enumitem}
\usepackage{listings}
\usepackage{fancyhdr}
\usepackage{setspace}
\usepackage{titlesec}
\usepackage{hyperref}
\usepackage{xurl}
\usepackage[numbers,sort&compress]{natbib}
\usepackage{cleveref}

% ---------- 全局样式 ----------
\definecolor{accent}{RGB}{31,78,121}
\definecolor{draftred}{RGB}{180,20,20}
\definecolor{codegray}{RGB}{246,247,249}
\hypersetup{hidelinks,pdfauthor={},pdftitle={海上油田人员直升机运载计划编排}}
\setstretch{1.25}
\setlength{\parindent}{2em}
\setlength{\parskip}{0pt}
\setlist{nosep,leftmargin=2em}
\captionsetup{font=small,labelsep=quad}
\renewcommand{\arraystretch}{1.18}

\pagestyle{fancy}
\fancyhf{}
\fancyfoot[C]{\thepage}
\renewcommand{\headrulewidth}{0pt}
\renewcommand{\footrulewidth}{0pt}
\fancypagestyle{plain}{\fancyhf{}\fancyfoot[C]{\thepage}\renewcommand{\headrulewidth}{0pt}}

\titleformat{\section}{\centering\heiti\zihao{3}}{\thesection}{1em}{}
\titleformat{\subsection}{\heiti\zihao{4}}{\thesubsection}{0.8em}{}
\titleformat{\subsubsection}{\heiti\zihao{-4}}{\thesubsubsection}{0.8em}{}
\titlespacing*{\section}{0pt}{1.4ex plus .4ex}{1.0ex}
\titlespacing*{\subsection}{0pt}{1.2ex plus .3ex}{0.7ex}

\newcommand{\papertitle}{海上油田人员直升机运载计划编排的分层建模与优化}

% 图片不存在时仍可编译；正式提交前必须替换全部占位框。
\newcommand{\figplaceholder}[2][0.82\linewidth]{%
  \fbox{\begin{minipage}[c][4.2cm][c]{#1}\centering
  \textcolor{gray}{#2\\（替换为自动生成的 PDF/PNG 图片）}
  \end{minipage}}}

% 附录代码：缺失时给出显眼提示，而不是中断论文编译。
\newcommand{\inputcode}[2]{%
  \IfFileExists{#2}{\lstinputlisting[caption={#1}]{#2}}%
  {\par\noindent\textcolor{draftred}{【缺少代码文件：\texttt{\detokenize{#2}}】}\par}}

\lstset{
  language=Python,
  basicstyle=\ttfamily\zihao{-4},
  keywordstyle=\color{accent},
  commentstyle=\color{gray},
  stringstyle=\color{teal},
  backgroundcolor=\color{codegray},
  frame=single,
  framerule=0.2pt,
  breaklines=true,
  showstringspaces=false,
  numbers=left,
  numberstyle=\tiny\color{gray},
  xleftmargin=2em,
  framexleftmargin=1.5em
}

\crefname{figure}{图}{图}
\crefname{table}{表}{表}
\crefname{equation}{式}{式}
\renewcommand{\figurename}{图}
\renewcommand{\tablename}{表}
\renewcommand{\refname}{参考文献}

\begin{document}
\pagenumbering{arabic}
\setcounter{page}{1}

% ================================================================
% 第1页：摘要专用页。标题、摘要、关键词总计不得超过一页。
% 不得出现姓名、学校、指导教师等身份信息。
% ================================================================
\thispagestyle{plain}
\begin{center}
  {\heiti\zihao{2}\papertitle\par}
  \vspace{1.2em}
  {\heiti\zihao{3}摘\quad 要}
\end{center}

本文针对海上油田人员出海、海返与设施间穿梭的复杂联合运输问题，建立了一个“需求聚合--路线构造--机型特点与加油联合判定--有限机队排班--独立校验”的统一求解框架。按照题目要求，模型以总飞机使用时间为首要目标，按词典序进一步优化人员总在途时间、总架次数和总燃油消耗量，并以座位利用率作为前述指标相同时的方案评价量；同时设计独立仿真校验器逐航段复查容量、油量、停靠、取送顺序与时间窗约束。

对于问题一，先以单设施容量动态规划和广义节约算法构造可行初解，再以关联度引导ALNS跨越局部盆地；随后收集多轨迹精英路线，求解带人员分配模式的精确路线池主问题，并用标准ALNS对主问题解作词典序精修。最终方案覆盖全部1600人且零违规，总飞机使用时间为14730 min、人员总在途时间为121363 min、共89个架次、燃油消耗118624.4 kg，座位利用率为49.430\%。相对广义节约初解，首要目标下降6.435\%。有限路线池LP值为14208.637 min，对应3.669\%的池内参考间隙，但不作为原问题全局最优性证明。

对于问题二，在候选路线整数主问题的基础上，采用4路线精确局部MILP修复、几何排序的1--5设施有界候选、质量精英重组、全局最优重启、停滞检查点续算及低频5路线强化。第三轮中普通低利用率4路线修复完成$4\to3$吸收，将架次数由96降至95；后续定向5路线与几何续算继续改善主、次目标。最终以95个架次完成全部4000人运输，总飞机使用时间为17218 min、人员总在途时间为254656 min、燃油消耗132473.9 kg，座位利用率为92.703\%，最终CSV独立校验零违规。相对三类需求分开运输基线，飞机时间下降61.031\%；相对第二轮17595 min再下降2.143\%。所有求解界仅属于有限候选主问题或当次局部邻域，不作为原问题全局最优性证明。

对于问题三，将问题分解为候选架次生成和具体飞机排班两层，并建立“第一阶段--第二阶段--删除临时人员投影--第一阶段再优化”的反馈闭环。程序采用12启动Top-$3$构造、时间感知排班与同日$k\to m$结构重构；随后以精英池邻域搜索、异构参数筛选（20组配置$\times$2种子）、4岛深度自适应大邻域、Optional Rescue与强制P0反馈继续深挖，其中跨日结构搜索贡献了深度阶段几乎全部全局改进。第一阶段最终以162个架次完成3840名非临时人员，基准$T_0=28728$ min，较旧深度解30546 min累计减少1818 min；第二阶段在同一上界下服务157名临时人员，飞机时间仍为28728 min。机场--机型分层连续多商品流给出全局有效下界14125 min，对应认证区间50.832\%；157人为最终固定架次结构下的可证明最优，不限定结构的158人及以上可行性仍未关闭。两阶段输出均通过独立硬约束校验。

\vspace{0.8em}
\noindent\textbf{关键词：}异构车辆路径；取送问题；时间窗；自适应大邻域搜索；区间排班；整数规划

\clearpage

% ================================================================
% 正文从第2页开始。官方要求：不要目录，正文控制在30页以内。
% ================================================================
\section{问题重述与分析}

\subsection{问题背景与任务}
题目给出 3 座陆地机场、52 座海上设施、3 种直升机和人员运输需求与任务，三个问题要求依次完成单向出海运输、出海--海返--穿梭联合运输，以及带时间窗和有限机队约束的多日排班。本文不重复题面叙述，而将三问概括为逐层增加约束的组合优化问题。

\begin{table}[htbp]
  \centering
  \caption{三个问题的递进关系}
  \label{tab:question-relation}
  \begin{tabularx}{\linewidth}{c X X}
    \toprule
    问题 & 核心决策 & 新增约束与目标 \\
    \midrule
    Q1 & 机场、机型、路线、乘客批次与加油选择 & 单向配送；容量、油量、最多 5 次海上着陆\\
    Q2 & 联合取送路线与逐航段载荷 & 海返与穿梭；同架次取送先后；禁止换乘\\
    Q3 & 候选架次、起飞时刻和具体飞机 & 7 天时间窗、24 架飞机、营业时段、30 min 周转、临时任务可选 \\
    \bottomrule
  \end{tabularx}
\end{table}

\subsection{目标的优先关系}
分析题目可知，在问题一与问题二中，题目要求“在最小化总飞机使用时间的基础上”优化其他指标。因此本文使用具有明确优先级的词典序目标。对 Q1 和 Q2，依次比较
\begin{equation}
  \bigl(T_{\mathrm{air}},\ T_{\mathrm{passenger}},\ N_{\mathrm{flight}},\
  F_{\mathrm{fuel}},\ -U\bigr).
  \label{eq:lexicographic}
\end{equation}
其中Q1在邻域解之间直接比较完整目标向量；Q2首先单独求解一级目标，固定给定时限内得到的
最好一级整数值，再用由目标上界自动推导的支配权重编码人员在途时间、架次数和燃油的
先后顺序，而非使用人工调参权重。座位利用率不进入线性主问题，仅在前四项完全相同时
用于比较已还原的完整方案。
Q3 先保证所有非临时任务完成，再最小化其基准总使用时间 $T_0$；第二阶段约束 $T_{\mathrm{air}}\le T_0$，优先最大化已服务临时人员数。

\subsection{技术路线}
\begin{figure}[htbp]
  \centering
  \IfFileExists{figures/framework.pdf}
    {\includegraphics[width=0.26\linewidth]{figures/framework.pdf}}
    {\figplaceholder{统一求解框架：数据校验、需求聚合、可行解构造、ALNS、区间排班、整数重分配、Validator}}
  \caption{统一建模与求解框架}
  \label{fig:framework}
\end{figure}

\section{模型假设说明}
在不改变题面规则的前提下，作如下假设：
\begin{enumerate}[label=\arabic*.]
  \item 题目给出的两地点距离代表一次直接飞行的有效航程；只有路线中真实记录的停靠点才允许拆分航段。
  \item 飞行速度与单位距离油耗不随载客量、天气和时段变化，停靠期间除规定时长外不另计油耗。
  \item 同一停靠点严格按照“到达、检查安全余油、下客、上客、检查容量、加油或等待、离开”的顺序更新状态。
  \item LAND 表示陆地端机场可选；明确给出 A01、A02 或 A03 时不得更换机场。
  \item 所有时间均按题意逐航段向上取整至整数分钟；不在全程结束后统一取整。
  \item 不考虑机场、设施和加油站的并发容量限制，但 Q3 中同一架飞机的架次不可重叠，且须满足机场周转时间。
\end{enumerate}

其中“相关路线移除”不再使用固定机场罚项。设路线$r$的实际服务设施集合为$V(r)$，原始
距离相关度取
\begin{equation}
  d_{\mathrm{rel}}(r_i,r_j)
  =\min_{u\in V(r_i),\,v\in V(r_j)}d_{uv}.
  \label{eq:q1-related-distance}
\end{equation}
在Combined R3配置中，程序又以几何距离、容量余量、潜在路线消除、机场匹配和当前路线
状态构成Context V2分数，只在建立精确候选和MILP之前进行排序与预算裁剪。该分数不代替
航段物理、燃油标签或目标函数；被保留的候选仍须经过完全相同的精确修复和Validator。

\section{符号说明}

本文使用的主要集合、参数、决策变量及评价指标如
表~\ref{tab:symbols}所示。

\begin{longtable}{
    >{\centering\arraybackslash}p{0.20\linewidth}
    p{0.70\linewidth}
}
    \caption{主要符号及其含义}
    \label{tab:symbols}\\
    
    \toprule
    符号 & 含义 \\
    \midrule
    \endfirsthead
    
    \multicolumn{2}{c}{续表~\thetable\quad 主要符号及其含义}\\
    \toprule
    符号 & 含义 \\
    \midrule
    \endhead
    
    \midrule
    \multicolumn{2}{r}{续下页}
    \endfoot
    
    \bottomrule
    \endlastfoot

    \multicolumn{2}{l}{\textbf{集合与索引}} \\[2pt]

    $A$
    & 陆地机场集合，$A=\{\mathrm{A01},\mathrm{A02},\mathrm{A03}\}$。 \\

    $F$
    & 海上设施集合，$F=\{\mathrm{F001},\ldots,\mathrm{F052}\}$。 \\

    $R$
    & 可加油海上设施集合，且满足 $R\subseteq F$。 \\

    $V$
    & 全部实际地点的集合，$V=A\cup F$；其中
    \texttt{LAND} 表示任意陆地机场，不作为独立地点加入 $V$。 \\

    $K$
    & 可用直升机机型集合，$K=\{\mathrm{T1},\mathrm{T2},\mathrm{T3}\}$。 \\

    $H$
    & 问题三中可用的具体直升机集合。 \\

    $P$
    & 人员运输需求集合，$p\in P$ 表示一条人员运输需求。 \\

    $P_{\mathrm{temp}}$
    & 临时任务人员集合，且满足 $P_{\mathrm{temp}}\subseteq P$。 \\

    $\mathcal{R}_{pk}$
    & 人员 $p$ 使用机型 $k$ 时，满足燃油和停靠次数约束的
    候选运输路径集合。 \\

    $i,j$
    & 地点索引，$i,j\in V$。 \\

    $k$
    & 直升机机型索引，$k\in K$。 \\

    $u$
    & 具体直升机索引，$u\in H$。 \\

    $h$
    & 架次路线中的停靠位置序号。 \\

    $p$
    & 人员运输需求索引，$p\in P$。 \\

    $\mathcal{B}^{(1)},\mathcal{B}^{(2)}$
    & 分别为问题一和问题二的聚合需求批次集合，$b$ 为需求批次索引。 \\

    $\ell,\gamma$
    & 分别为ALNS算子索引和算子权重更新周期索引。 \\

    \midrule
    \multicolumn{2}{l}{\textbf{基础参数}} \\[2pt]

    $o_p,d_p$
    & 人员 $p$ 的起点和终点。 \\

    $e_p$
    & 人员 $p$ 所乘飞机离开其起点的最早允许时刻。 \\

    $l_p$
    & 人员 $p$ 抵达其终点的最晚允许时刻。 \\

    $c_p$
    & 人员 $p$ 的业务类型，取应急处置、增储上产、
    常规倒班或临时任务。 \\

    $\operatorname{prio}(c_p)$
    & 业务类型 $c_p$ 对应的优先等级，数值越小表示优先级越高。 \\

    $d_{ij}$
    & 地点 $i$ 到地点 $j$ 的题给直接飞行距离，单位为 km。 \\

    $s_k$
    & 机型 $k$ 的可用座位数。 \\

    $\nu_k$
    & 机型 $k$ 的巡航速度，单位为 km/h。 \\

    $\beta_k$
    & 机型 $k$ 的单位距离燃油消耗量，单位为 kg/km。 \\

    $G_k^{\max}$
    & 机型 $k$ 的油箱容量，单位为 kg。 \\

    $G_k^{\min}$
    & 机型 $k$ 到达任意地点时必须满足的最低安全余油量，
    单位为 kg。 \\

    $D_k^{\max}$
    & 机型 $k$ 满油出发时的安全可用航程，
    $D_k^{\max}=(G_k^{\max}-G_k^{\min})/\beta_k$，单位为 km。 \\

    $t_{ij}^{k}$
    & 机型 $k$ 从地点 $i$ 飞往地点 $j$ 的飞行时间，
    $t_{ij}^{k}=\left\lceil 60d_{ij}/\nu_k\right\rceil$，单位为 min。 \\

    $g_{ij}^{k}$
    & 机型 $k$ 从地点 $i$ 飞往地点 $j$ 的燃油消耗量，
    $g_{ij}^{k}=\beta_kd_{ij}$，单位为 kg。 \\

    $\underline{\tau}_i^k$
    & 计算乐观最短运输时间时，机型 $k$ 在地点 $i$ 的最短技术停靠
    时间，单位为 min。 \\

    $Q_b$
    & 需求批次 $b$ 包含的人员数量。 \\

    $A_b$
    & 需求批次 $b$ 允许使用的陆地机场集合；指定机场需求只能使用
    对应机场，\texttt{LAND} 需求满足 $A_b=A$。 \\

    $f_b$
    & 问题一中需求批次 $b$ 的目标海上设施。 \\

    $a_b$
    & 问题一中与需求批次 $b$ 绑定的指定陆地机场。 \\

    $o_b,d_b$
    & 问题二中需求批次 $b$ 的起点和终点。 \\

    \midrule
    \multicolumn{2}{l}{\textbf{需求聚合与时间窗指标}} \\[2pt]

    $Q_{ij}$
    & 起点为 $i$、终点为 $j$ 的聚合人员需求数量。 \\

    $W_p$
    & 人员 $p$ 的时间窗宽度，
    $W_p=l_p-e_p$，单位为 min。 \\

    $\underline{T}_p$
    & 人员 $p$ 在放松飞机调位、返场和资源竞争等约束后的
    乐观最短运输时间，单位为 min。 \\

    $S_p$
    & 人员 $p$ 的乐观时间松弛量，
    $S_p=W_p-\underline{T}_p$，单位为 min。 \\

    $\Pi_p$
    & 人员 $p$ 的任务排序指标，
    $\Pi_p=(\operatorname{prio}(c_p),l_p,S_p)$，
    按升序进行词典序比较。 \\

    \midrule
    \multicolumn{2}{l}{\textbf{路线状态变量}} \\[2pt]

    $r=(v_{r0},\ldots,v_{r,m_r+1})$
    & 架次 $r$ 的有序停靠节点序列，其中
    $v_{r0}=v_{r,m_r+1}=a(r)\in A$，中间节点均属于 $F$。 \\

    $k(r)$
    & 执行架次 $r$ 的直升机机型。 \\

    $a(r)$
    & 架次 $r$ 的起降机场；架次从 $a(r)$ 起飞并返回同一机场。 \\

    $m_r$
    & 架次 $r$ 的海上着陆次数，满足 $m_r\leq5$。 \\

    $\eta_{br}$
    & 分配给架次 $r$ 运输的需求批次 $b$ 的人数，为非负整数。 \\

    $L_{rh}$
    & 架次 $r$ 离开第 $h$ 个停靠点之后的机上人数。 \\

    $\xi_{rh}$
    & 架次 $r$ 在第 $h$ 个海上停靠点的加油决策变量；加油时取1，
    否则取0。 \\

    $\tau_{rh}$
    & 架次 $r$ 在第 $h$ 个停靠点的实际停靠时间，单位为 min。 \\

    $T_{br}^{\mathrm{pass}}$
    & 批次层单人在途时间；用于问题一、二的聚合需求评价，单位为 min。 \\

    $T_r^{\mathrm{air}}$
    & 架次 $r$ 的飞机使用时间，即从机场起飞至返回机场的总时长，
    包括飞行、海上停靠、加油和必要等待时间，单位为 min。 \\

    $G_{rh}^{\mathrm{arr}}$
    & 架次 $r$ 到达第 $h$ 个停靠点时的剩余燃油量，单位为 kg。 \\

    $G_{rh}^{\mathrm{dep}}$
    & 架次 $r$ 离开第 $h$ 个停靠点时的剩余燃油量，单位为 kg。 \\
\midrule
\multicolumn{2}{l}{\textbf{问题二决策变量与状态变量}} \\[2pt]

$\pi_{br}$
& 需求批次 $b$ 在架次 $r$ 上的上机位置序号。 \\

$\delta_{br}$
& 需求批次 $b$ 在架次 $r$ 上的下机位置序号；
当 $\eta_{br}>0$ 时满足
$\pi_{br}<\delta_{br}$。 \\

$N_{rh}^{\mathrm{on}}$
& 架次 $r$ 在第 $h$ 个停靠点的上机人数，即
\[
N_{rh}^{\mathrm{on}}=\sum_{b:\,\pi_{br}=h}\eta_{br}.
\] \\

$N_{rh}^{\mathrm{off}}$
& 架次 $r$ 在第 $h$ 个停靠点的下机人数，即
\[
N_{rh}^{\mathrm{off}}=\sum_{b:\,\delta_{br}=h}\eta_{br}.
\] \\

$\mathcal C$
& 问题二经燃油标签搜索后得到的候选闭合路线集合，$c\in\mathcal C$。 \\

$y_c$
& 问题二候选路线 $c$ 的使用架次数，为非负整数。 \\

$x_{bc}$
& 问题二批次 $b$ 分配至候选路线 $c$ 的人数，为非负整数。 \\

$\kappa_{bch}$
& 批次 $b$ 在候选路线 $c$ 的第 $h$ 个航段上占用座位时取1，否则取0。 \\

$\theta_{bc}$
& 批次 $b$ 采用候选路线 $c$ 时的单人在途时间，单位为 min。 \\

$D,\mathcal B_D,\mathcal C_D$
& 问题二一次局部搜索中被破坏的路线集合、对应局部需求集合和有界候选路线集合。 \\

\midrule
\multicolumn{2}{l}{\textbf{问题一辅助函数与ALNS算法变量}} \\[2pt]

$\lambda$
& 容量动态规划中新增架次的载客人数，满足
$1\leq\lambda\leq\min\{n,s_k\}$。 \\

$\omega_{af}$
& 起点为\texttt{LAND}、终点为设施 $f$ 的人员中，分配至机场 $a$ 的人数。 \\

$\Phi_{af}(n)$
& 从机场 $a$ 向设施 $f$ 运输剩余 $n$ 人时，容量动态规划的最优
词典序目标向量。 \\

$w_{\ell}^{(\gamma)}$
& 第 $\gamma$ 个统计周期内第 $\ell$ 个破坏或修复算子的自适应权重。 \\

$\rho$
& ALNS 算子权重的更新系数，满足 $0<\rho\leq 1$。 \\

$\sigma_{\ell}^{(\gamma)}$
& 第 $\gamma$ 个统计周期内第 $\ell$ 个算子获得的累计奖励分数。 \\

$n_{\ell}^{(\gamma)}$
& 第 $\gamma$ 个统计周期内第 $\ell$ 个算子被调用的次数。 \\[2pt]

$\Theta,\alpha$
& 分别为模拟退火的当前温度和降温系数，温度按照
$\Theta\leftarrow\alpha\Theta$ 更新。 \\[2pt]

$S_{\mathrm{init}}^{(1)}$
& 问题一经广义节约合并得到、供 ALNS 继续改进的初始可行方案。 \\

$V(r),d_{\mathrm{rel}}(r_i,r_j)$
& 路线$r$的实际服务设施集合，以及两条路线之间基于原始设施距离的相关度距离。 \\

$\mathcal G_D,\mathcal R_D^{\mathrm{cand}}$
& 一次ALNS破坏后形成的待修复需求组集合与候选路线集合。 \\

$Q_g^D$
& 待修复需求组 $g$ 的人数。 \\

$\zeta_c$
& 整数规划修复中候选路线 $c$ 的使用架次数。 \\

$\chi_{gc}$
& 整数规划修复中需求组 $g$ 分配至候选路线 $c$ 的人数。 \\

$\vartheta_{cg}$
& 候选路线 $c$ 运送需求组 $g$ 时的单人在途时间，单位为 min。 \\

$\mathcal R^{\mathrm{pool}},\mathcal P$
& 分别为问题一跨轨迹去重后的精英物理路线集合和带人员分配的路线模式集合。 \\

$z_p$
& 精确路线池主问题中是否选用分配模式 $p$ 的0--1变量。 \\

$a_{bp}$
& 路线模式 $p$ 对问题一需求批次 $b$ 的覆盖人数。 \\

\midrule
\multicolumn{2}{l}{\textbf{问题二结果比较指标}} \\[2pt]

$T_{\mathrm{separate}}$
& 出海、海返和穿梭三类需求分开运输时的总飞机使用时间。 \\

$T_{\mathrm{joint}}$
& 采用联合优化算法时的总飞机使用时间。 \\

$I_{\mathrm{joint}}$
& 联合运输相对分开运输的飞机使用时间降低率，
\[
I_{\mathrm{joint}}
=
\frac{T_{\mathrm{separate}}-T_{\mathrm{joint}}}
     {T_{\mathrm{separate}}}
\times 100\%.
\] \\
\midrule
\multicolumn{2}{l}{\textbf{问题三新增集合、参数与决策变量}} \\[2pt]

$\mathcal{D}$
& 问题三的规划日集合，$\mathcal{D}=\{0,1,\ldots,6\}$，分别对应
2026年8月3日至8月9日。 \\

$q$
& 规划日索引，$q\in\mathcal{D}$。 \\

$\Omega_q$
& 第 $q$ 日相对规划起点的分钟偏移，$\Omega_q=1440q$。 \\

$P_{\mathrm{req}}$
& 必须完成运输的非临时人员集合，
$P_{\mathrm{req}}=P\setminus P_{\mathrm{temp}}$。 \\

$\mathcal{C}$
& 路线层生成的时间可行候选架次集合。 \\

$P_r$
& 候选架次 $r$ 计划运输的人员集合。 \\

$\mathcal{C}_p$
& 能够运输人员 $p$ 的候选架次集合，
$\mathcal{C}_p=\{r\in\mathcal{C}:p\in P_r\}$。 \\

$q(r)$
& 候选架次 $r$ 所属的规划日。 \\

$a(u),k(u)$
& 具体飞机 $u$ 的所属机场和机型。 \\

$H_r$
& 能够执行候选架次 $r$ 的具体飞机集合，
$H_r=\{u\in H:a(u)=a(r),\,k(u)=k(r)\}$。 \\

$x_r$
& 候选架次选择变量；若选择候选架次 $r$，则 $x_r=1$，
否则 $x_r=0$。 \\

$y_{ru}$
& 架次--飞机分配变量；若候选架次 $r$ 由具体飞机 $u$ 执行，
则 $y_{ru}=1$，否则 $y_{ru}=0$。 \\

$T_{rh}^{\mathrm{arr}}$
& 架次 $r$ 到达第 $h$ 个停靠点的绝对分钟时刻，以
2026年8月3日00:00为时间原点。 \\

$T_{rh}^{\mathrm{dep}}$
& 架次 $r$ 离开第 $h$ 个停靠点的绝对分钟时刻。 \\

$\pi_{pr}$
& 问题三中人员 $p$ 在候选架次 $r$ 上的上机位置序号；它是
批次层符号 $\pi_{br}$ 的个人层对应量。 \\

$\delta_{pr}$
& 问题三中人员 $p$ 在候选架次 $r$ 上的下机位置序号；它是
批次层符号 $\delta_{br}$ 的个人层对应量，满足 $\pi_{pr}<\delta_{pr}$。 \\

$T_{pr}^{\mathrm{pass}}$
& 问题三个人层实际在途时间；它与批次层
$T_{br}^{\mathrm{pass}}$ 含义相同但索引粒度不同，
$T_{pr}^{\mathrm{pass}}=
T_{r,\delta_{pr}}^{\mathrm{arr}}-T_{r,\pi_{pr}}^{\mathrm{dep}}$。 \\

$I_{ru}$
& 候选架次 $r$ 分配给飞机 $u$ 时建立的可选时间区间；区间末端
包含30 min机场周转时间，其存在变量为 $y_{ru}$。 \\

$z_p$
& 临时人员服务变量；若临时人员 $p$ 被纳入最终运输方案，
则 $z_p=1$，否则 $z_p=0$。 \\

$w_{pr}$
& 固定架次人员重分配变量；若人员$p$分配到架次$r$，则取1，否则取0。 \\

$\mathcal{P}_{\mathrm{req}}(S)$
& 将方案$S$中的临时人员分配全部删除、同时保持航班结构和时刻不变的
非临时任务投影算子。 \\

$S_{\mathrm{base}}$
& 第一阶段完成全部非临时任务后得到的最佳可行基准方案。 \\

$T_0$
& 基准方案 $S_{\mathrm{base}}$ 的总飞机使用时间，
$T_0=T_{\mathrm{air}}(S_{\mathrm{base}})$，单位为 min。 \\

$M_p$
& 最终排班中已服务人员 $p$ 的实际送达裕度；与候选生成前计算的
乐观松弛量 $S_p$ 不同，
$M_p=l_p-T_{r,\delta_{pr}}^{\mathrm{arr}}$，单位为 min。 \\

$R_{\mathrm{temp}}$
& 临时任务满足率，
$R_{\mathrm{temp}}=
\left(\sum_{p\in P_{\mathrm{temp}}}z_p/160\right)\times100\%$。 \\

$\mathcal{L}=A\times K$
& 增强下界模型中的机场--机型飞机流层集合；层$\ell=(a,k)$仅允许从机场$a$
出发并在$a$与海上设施之间流动。 \\

$\mathcal{E}_{\ell}$
& 飞机流层$\ell$中的有向飞行弧集合。 \\

$c_{ij}^{\ell}$
& 飞机流层$\ell=(a,k)$在弧$(i,j)$上的下界成本，为逐航段整数飞行时间与
到达海上设施的最小停靠时间之和，单位为 min。 \\

$n_{ak}$
& 机场$a$所属机型$k$的可用飞机数量。 \\

$X_{ij}^{\ell}$
& 增强下界模型中层$\ell$在有向弧$(i,j)$上的连续飞机流量。 \\

$F_{ij}^{d}$
& 以节点$d$为共同终点的聚合人员在弧$(i,j)$上的连续人员流量。 \\

$b_i^d$
& 终点聚合商品$d$在节点$i$的净供给量；起点为正、共同终点为负，
其余节点为0。 \\

$LB_3^{\mathrm{MCF}}$
& 机场--机型分层连续多商品流模型给出的原问题第一阶段全局下界。 \\

$n_r^{\mathrm{LP}}$
& 有限候选路线池参考模型中路线$r$的连续执行次数。 \\

$w_{gr}^{\mathrm{LP}}$
& 有限候选路线池参考模型中OD组$g$分配给路线$r$的连续人数。 \\

    \midrule
    \multicolumn{2}{l}{\textbf{评价指标}} \\[2pt]

    $T_{\mathrm{air}}$
    & 全部架次的总飞机使用时间，单位为 min。 \\

    $T_{\mathrm{passenger}}$
    & 全部已运输人员的总在途时间，单位为 min。 \\

    $N_{\mathrm{flight}}$
    & 最终运输方案中的总架次数。 \\

    $F_{\mathrm{fuel}}$
    & 全部架次的总燃油消耗量，单位为 kg。 \\

    $U$
    & 全部架次的座位利用率，即总客座公里与总可用座公里之比。 \\

\end{longtable}

为避免混淆，本文将 $q$ 专用于问题三的规划日索引；问题一模拟退火接受准则中
“目标向量首个发生劣化的分量”统一记为 $j^\star$。此外，带下标 $b$ 的
$\pi_{br},\delta_{br},T_{br}^{\mathrm{pass}}$ 是问题一、二的批次层量，带下标
$p$ 的 $\pi_{pr},\delta_{pr},T_{pr}^{\mathrm{pass}}$ 是问题三的个人层量，
二者物理含义相同但建模粒度不同。

\section{数据预处理与探索性分析}

\subsection{数据预处理}
首先，本文建立统一参数文件
\texttt{configs/problem.json}，集中保存题目给出的机场、设施、
加油点、机型参数、停靠规则及问题三的机队和运营时段等信息，
避免不同程序重复硬编码参数而产生口径差异。

随后，由 \texttt{data\_pipeline.py} 建立受保护的原始数据副本，
记录4个原始CSV文件的大小和SHA-256摘要。程序重复运行时会重新
核对文件摘要，若原始数据发生非预期修改则立即报错。
\subsection{数据一致性校验}

接着，我们验证了距离矩阵节点、对称性、对角零和缺失值；随后检查人员编号唯一性、起终点合法性、时间窗顺序及 Q1--Q3 之间的人员对应关系。计算过程中始终使用题给相邻停靠点距离，不对距离矩阵作最短路闭包，结果如表4所示。所有校验均通过，说明题目数据完整且一致。

\begin{table}[htbp]
  \centering
  \caption{数据质量检验结果}
  \label{tab:data-process-true}
  \begin{tabular}{lc}
    \toprule
    项目 & 检验结果  \\
    \midrule
    距离矩阵 & $55\times55$ \\
    节点 & 3 个机场、52 个设施  \\
    缺失值 & 0  \\
    负距离 & 0  \\
    对角线 & 全为0  \\
    对称性 & 满足  \\
    非对角距离 & 14～579 km  \\
    Q1 人数 & 1600  \\
    Q2/Q3 人数 & 4000  \\
    人员编号重复 & 0  \\
    非法起终点 & 0  \\
    Q3 非法时间窗 & 0  \\

    \bottomrule
  \end{tabular}
\end{table}

我们还发现距离矩阵不完全满足三角不等式：
147 个无序节点对存在经第三点更短的情况；
对应 294 个有向节点对；
共 442 个具体违例三元组。但是由于如果飞机经过中间设施，该设施必须真实写入路线，并计入着陆、停靠、油量和时间。故没有将距离改成最短距离。
\subsection{需求结构与 OD 聚合}

按照人员起终点的节点类型，本文将运输需求划分为出海、海返和穿梭三类。设人员 $p$ 的起点和终点分别为 $o_p$ 和 $d_p$，机场集合和海上设施集合分别为 $A$ 和 $F$，则人员运输类型定义为

\begin{equation}
\operatorname{type}(p)=
\begin{cases}
\mathrm{outbound},
& o_p\in A\cup\{\mathrm{LAND}\},\quad d_p\in F,\\
\mathrm{inbound},
& o_p\in F,\quad d_p\in A\cup\{\mathrm{LAND}\},\\
\mathrm{shuttle},
& o_p\in F,\quad d_p\in F.
\end{cases}
\label{eq:demand-type}
\end{equation}
设 $Q_{ij}$ 表示从节点 $i$ 到节点 $j$ 的人员数量，则
\begin{equation}
Q_{ij}
=
\sum_{p\in P}
\mathbb{I}(o_p=i,\ d_p=j),
\label{eq:od-aggregation}
\end{equation}
其中，$\mathbb{I}(\cdot)$ 为示性函数。


问题一包含 1600 名出海人员；问题二和问题三均包含 4000 名人员，其中出海、海返和穿梭人员分别为 1600、1600 和 800 人，占比分别为 $40\%$、$40\%$ 和 $20\%$。各类运输需求及其唯一 OD 数量如表~\ref{tab:demand-od-structure} 所示。

\begin{table}[htbp]
    \centering
    \caption{不同运输类型的人员数量与唯一 OD 数量}
    \label{tab:demand-od-structure}
    \begin{tabular}{ccccc}
        \toprule
        问题 & 运输类型 & 人员数量 & 人数占比  & 唯一 OD 数量 \\
        \midrule
        Q1 & 出海 & 1600 & $100\%$  & 104 \\
        \midrule
        \multirow{3}{*}{Q2}
            & 出海 & 1600 & $40\%$  & 104 \\
            & 海返 & 1600 & $40\%$  & 104 \\
            & 穿梭 & 800 & $20\%$  & 56 \\
        \midrule
        Q2 合计 & -- & 4000 & $100\%$ & 264 \\
        \bottomrule
    \end{tabular}
\end{table}
值得注意的是，Q3虽然具有相同的264个OD，但同一OD内不同人员的时间窗和业务类型可能不同，因此Q3仅使用OD作为检索索引，不直接合并为不可区分的批量需求。

问题二中的出海和海返需求规模相等，说明同一架次在海上设施送下出海人员后，可以利用刚释放的座位接运海返人员；穿梭人员则可进一步利用设施之间航段的剩余运力。因此，相比将三类需求分别规划，联合运输具有降低空载航段和提高座位利用率的潜力。

\begin{figure}[htbp]
    \centering
    \includegraphics[width=0.6\textwidth]
    {figures/q2_direction_counts.png}
    \caption{问题二不同运输方向的人员需求量}
    \label{fig:q2-direction-counts}
\end{figure}

由图~\ref{fig:q2-direction-counts} 可见，出海与海返需求在总量上具有良好的对称性，为构造“出海送人、海上接人、返程载客”的混合架次提供了客流基础。但这种总量对称并不意味着各设施的人员流量完全平衡，因此还需进一步分析各 OD 和设施的需求分布。


\subsubsection{OD 需求聚合方法}

通过分析表5，我们发现
Q1 前 10 个 OD 占 23.5\%，Q2 前 10 个 OD 仅占 10.5\%。说明需求具有一定重复性，可以按 OD 聚合降低搜索规模，但不能只针对少数热门 OD 规划。
因此，本文在搜索层按照起终点对聚合人员数量，以减少重复计算；在最终方案生成阶段再将批量需求还原为人员编号，保证每名人员均具有唯一、可追踪的运输安排。
\subsection{LAND 端点与机场选择分析}
对于起点或终点为 LAND 的人员，本文保留 A01、A02 和 A03 三个候选机场，不在数据预处理阶段固定实际机场。最近机场仅作为路线构造的启发式特征。但是若将Q1的1353名灵活出海人员全部分配至最近机场，则 A01、A02 和 A03 分别承担约865、193和295人，A01的需求明显偏高。由此可见，机械采用最近机场规则可能造成机场间任务分布不均，因此后续将 LAND 对应机场与路线、机型共同优化。
\subsection{航段特征与燃油可达性}
对每个机型和有序地点对预计算飞行时间（向上取整）与燃油消耗：
\begin{equation}
  t_{ij}^{k}=\left\lceil\frac{60d_{ij}}{\nu_k}\right\rceil,
  \qquad g_{ij}^{k}=\beta_kd_{ij}.
  \label{eq:precompute}
\end{equation}
因为有 55 个节点、每个节点有 54 个不同目的地、3 种机型，所以共生成：
\begin{equation}
  55\times54\times3=8910
\end{equation}
条航段特征。

同时，机型不能把油箱里的燃油全部消耗完，因为到达每一个地点时都必须保留最低安全余油。
因此，满油状态下的安全可用航程为：
\begin{equation}
  D_k^{\max}=\frac{G_k^{\max}-G_k^{\min}}{\beta_k}.
\end{equation}
计算得到T1、T2和T3的安全可用航程分别为250.00 km、400.00 km和482.76 km。进一步统计三种机型从不同机场出发时可以直接闭合往返的设施数量，结果如图所示。
\begin{figure}[htbp]
    \centering
    \includegraphics[width=0.66\textwidth]
    {figures/fuel_direct_roundtrip_coverage.png}
    \caption{不同机场和机型的直接闭合往返覆盖设施数}
    \label{fuel_direct_roundtrip_coverage}
\end{figure}
T1虽然具有最高飞行速度，但由于安全可用航程最短，从三座机场出发均不能在不加油的条件下直接闭合往返任一设施。T2从A01、A02和A03出发分别可以直接往返13、7和0座设施，T3则分别可以直接往返32、28和14座设施。这说明机型的速度优势、容量优势和燃油可达性之间存在明显权衡，不能在路线生成后仅根据飞行速度被动指定机型。

对于无法直接闭合往返的设施，本文进一步在最多5次海上着陆范围内搜索包含技术停靠和加油的闭合路线。搜索状态记录当前位置、剩余燃油、目标设施访问状态、加油状态和已发生的海上着陆次数。
飞机到达下一地点时的剩余燃油满足
\begin{equation}
G_{r,h+1}^{\mathrm{arr}}
=
G_{rh}^{\mathrm{dep}}
-\beta_{k(r)}d_{v_{rh},v_{r,h+1}}
\ge G_{k(r)}^{\min}.
\end{equation}

若下一节点为可加油设施且选择加油，则：
\begin{equation}
  G_{j}^{k,\mathrm{dep}}=G_k^{\max}.
\end{equation}
否则：
\begin{equation}
  G_{j}^{k,\mathrm{dep}}=G_{j}^{k,\mathrm{arr}}.
\end{equation}
对于机场 $a$、设施 $f$ 和机型 $k$，若不经过其他海上设施且
不进行中途加油，其直接闭合往返可行条件为
\begin{equation}
2\beta_kd_{af}\le G_k^{\max}-G_k^{\min},
\end{equation}
等价地，
\begin{equation}
2d_{af}\le D_k^{\max}.
\end{equation}
结果表明，在允许使用8座指定加油设施的条件下，9种机场—机型组合均能在不超过5次海上着陆的范围内访问全部52座设施并闭合返场。但各机型所需的技术停靠次数差异较大：T1访问部分远程设施需要3～5次海上着陆，T2对全部设施最多需要3次，T3则最多需要2次。
上述结果说明，燃油可达性与停靠次数会显著影响机型的实际使用时间。T1虽然飞行速度较高，但可能因额外技术停靠和加油而失去时间优势；T3虽然飞行速度较低，却可能通过减少技术停靠获得更优的总架次时间。因此，后续路线优化应联合决定起降机场、访问顺序、机型和加油选择，而不能将机型选择视为路线生成后的独立步骤。

需要说明的是，本节得到的闭合路线可达性只表示燃油层面存在可行访问方式，尚未考虑人员容量、取送顺序、时间窗和总目标值，不能直接作为最终运输路线。
\subsection{Q3时间窗与任务紧迫性分析}
问题三在人员运输需求中增加了任务类型、最早登机时间和最晚到达时间等约束。
与前两问相比，问题三不仅需要确定人员的运输路径，还需要保证每项任务在规定的
时间窗内完成。

\subsubsection{时间窗参数定义}

设问题三中的人员运输任务集合为 $P$。对于任意任务
$p\in P$，记其属性为
\begin{equation}
    p=(o_p,d_p,e_p,l_p,c_p),
\end{equation}
其中，$o_p$ 和 $d_p$ 分别表示任务 $p$ 的出发设施和目标设施，
$e_p$ 表示最早允许登机时间，$l_p$ 表示最晚允许到达时间，
$c_p$ 表示任务类型。

定义任务 $p$ 的时间窗宽度为
\begin{equation}
    W_p=l_p-e_p.
    \label{eq:window_width}
\end{equation}
时间窗宽度反映了任务可用于运输的总时间。一般而言，$W_p$ 越小，
任务可以调整的空间越有限，其调度紧迫程度越高。

但是，仅根据时间窗宽度不能充分判断任务是否紧迫。例如，两个任务可能具有相同的
时间窗宽度，但由于运输距离、机型性能和途中加油需求不同，完成运输所需的时间可能
存在显著差异。因此，本文进一步引入乐观最短运输时间和时间松弛量。

\subsubsection{乐观最短运输时间}

对于任务 $p$，定义其乐观最短运输时间为 $\underline{T}_p$。
该时间表示在暂不考虑其他任务竞争、机队数量限制和航班周转冲突的条件下，
将人员从 $o_p$ 运送到 $d_p$ 所需的最短时间。

在计算 $\underline{T}_p$ 时，综合考虑以下因素：

\begin{enumerate}
    \item 不同机型的巡航速度和航段飞行时间；
    \item 各机型的燃油容量、单位距离耗油量和安全剩余油量；
    \item 允许使用的陆地机场；
    \item 海上加油设施及相应的加油时间；
    \item 中途技术经停次数不超过题目规定上限；
    \item 每一航段在燃油约束下均具有可达性。
\end{enumerate}

对于采用机型 $k$、经过节点序列
\[
    r=(v_0,v_1,\ldots,v_m),
\]
且满足 $v_0=o_p$、$v_m=d_p$ 的运输路径，其总时间可表示为
\begin{equation}
    T_{pkr}
    =
    \sum_{h=0}^{m-1}
    \left\lceil
    \frac{60d_{v_hv_{h+1}}}{\nu_k}
    \right\rceil
    +
    \sum_{h=1}^{m-1}\underline{\tau}_{v_h}^{k},
    \label{eq:q3_path_time}
\end{equation}
其中，$d_{v_hv_{h+1}}$ 为相邻节点之间的航程，$\nu_k$ 为机型 $k$
的巡航速度，$\underline{\tau}_{v_h}^{k}$ 表示机型 $k$ 在节点 $v_h$ 发生技术经停时
所需要的额外时间，包括起降、人员上下机和加油等作业时间。

在所有满足燃油约束和经停次数约束的机型与路径中取最小值，得到
\begin{equation}
    \underline{T}_p
    =
    \min_{k\in K},\,r\in\mathcal{R}_{pk}
    T_{pkr},
    \label{eq:optimistic_time}
\end{equation}
其中，$ K $ 为可用机型集合，
$\mathcal{R}_{pk}$ 为任务 $p$ 使用机型 $k$ 时的可行运输路径集合。

需要指出的是，$\underline{T}_p$ 是一种乐观下界。其计算暂未考虑同一架直升机
连续执行多个任务时的衔接过程，也未考虑机队规模、人员容量、航班周转时间和
其他任务对资源的竞争。为保证该指标确实构成乐观下界，计算时暂不计入飞机到达人员起点
之前的调位时间，也不计入人员送达后的返场时间；对于海上起点，
将起点出发油量放松为有利状态。因此，该指标可用于任务可行性预筛选和紧迫程度评价，
但不能代替最终调度模型的可行性检验。$\underline{T}_p$ 仅用于
衡量人员运输任务本身的最低时间需求，不能作为完整架次的可行时间。
\subsubsection{任务时间松弛量}

根据时间窗宽度和乐观最短运输时间，定义任务 $p$ 的时间松弛量为
\begin{equation}
    S_p=W_p-\underline{T}_p.
    \label{eq:q3_slack}
\end{equation}

时间松弛量表示在最理想运输条件下，任务完成后仍然剩余的可调度时间。
根据 $S_p$ 的取值，可以得到以下判断：

\begin{equation}
    \begin{cases}
        S_p<0, & \text{任务在当前运输条件下不具备基本可行性};\\
        S_p\approx 0, & \text{任务几乎没有调整空间，紧迫程度较高};\\
        S_p>0, & \text{任务具有一定的时间调整空间}.
    \end{cases}
\end{equation}

对问题三的全部人员任务进行统计，得到不同业务类型的时间窗和松弛量特征，
如表~\ref{tab:q3_time_window_summary} 所示。

\begin{table}[htbp]
    \centering
    \small
    \caption{问题三不同业务类型的时间窗与松弛量统计}
    \label{tab:q3_time_window_summary}
    \begin{tabular}{lrrrrr}
        \toprule
        任务类型
        & 人数
        & 时间窗中位数/min
        & 松弛量中位数/min
        & 最小松弛量/min
        & 松弛量小于60 min人数 \\
        \midrule
        应急处置 & 40   & 101.5 & 48.5  & 17   & 31 \\
        增储上产 & 360  & 181.0 & 129.0 & 65   & 0  \\
        常规倒班 & 3440 & 2922.0& 2883.0& 383  & 0  \\
        临时任务 & 160  & 4986.0& 4914.0& 1869 & 0  \\
        \bottomrule
    \end{tabular}
\end{table}

由表~\ref{tab:q3_time_window_summary} 可知，应急处置的时间窗中位数仅为
$101.5$ min，松弛量中位数为 $48.5$ min。在40项应急运输需求中，
有31项任务的松弛量不足60 min，其中3项任务的松弛量不足30 min。
因此，应急处置是问题三调度过程中最需要优先处理的任务类型。

增储上产的时间窗中位数为 $181$ min，松弛量中位数为 $129$ min，
其紧迫程度低于应急处置，但仍明显高于常规倒班和临时任务。
常规倒班与临时任务的时间窗口较宽，整体上具有较大的调整空间，
可以用于填补应急处置和增储上产完成后的剩余运力。

不同业务类型的松弛量分布如图~\ref{fig:q3_slack} 所示。

\begin{figure}[htbp]
    \centering
    \includegraphics[width=0.78\textwidth]
    {figures/q3_slack_by_task.png}
    \caption{不同业务类型的乐观时间松弛量分布}
    \label{fig:q3_slack}
\end{figure}

虽然全部任务的 $S_p$ 均不小于0，说明各项任务在忽略资源竞争时均具有基本
运输可行性，但这并不能保证所有任务可以同时完成。多个应急任务可能集中于相近
时间段，并竞争同一架直升机、同一机场或同一加油设施。因此，最终仍需在调度模型中
同时考虑机队数量、任务衔接和资源占用等约束。

\subsubsection{任务优先级排序规则}

综合考虑任务类型、最晚到达时间和时间松弛量，本文为每项任务构造如下排序指标：
\begin{equation}
    \Pi_p=
    \left(
        \operatorname{prio}(c_p),
        l_p,
        S_p
    \right).
    \label{eq:q3_priority}
\end{equation}
其中，$\operatorname{prio}(c_p)$ 表示任务类型的优先等级。
任务排序时依次按照以下规则进行：

\begin{enumerate}
    \item 优先安排业务等级更高的任务；
    \item 对于同一业务等级，优先安排最晚到达时间 $l_p$ 更早的任务；
    \item 若最晚到达时间相同，则优先安排松弛量 $S_p$ 更小的任务。
\end{enumerate}

按照题目给出的业务要求，任务类型优先级设置为
\begin{equation}
    \text{应急处置}
    \;>\;
    \text{增储上产}
    \;>\;
    \text{常规倒班}
    \;>\;
    \text{临时任务}.
\end{equation}

该排序规则将业务优先级、最早截止时间规则和最小松弛量规则相结合，
既能体现应急处置的优先性，又能降低时间窗较窄任务发生延误的概率，
可作为后续初始运输方案生成和邻域搜索中的任务选择依据。

\subsubsection{需求的日期分布}

进一步按照最早登机时间所在日期统计人员需求量，结果如
图~\ref{fig:q3_daily_pressure} 所示。8月3日的需求量为880人，
是整个规划期内需求最集中的日期；8月4日至8月8日的需求量分别为
537、571、554、581和583人，整体保持在较为稳定的水平；
8月9日的需求量为294人。

\begin{figure}[htbp]
    \centering
    \includegraphics[width=0.82\textwidth]
    {figures/q3_daily_pressure.png}
    \caption{问题三按最早登机日期统计的人员需求量}
    \label{fig:q3_daily_pressure}
\end{figure}

8月3日较高的需求量表明，规划期初可能面临更大的机队调度压力，
需要为该日期预留较多可用运力。不过，图~\ref{fig:q3_daily_pressure}
是按照任务的最早登机日期进行统计，而不是最终运输方案中的实际执行日期。
对于时间窗较宽的常规倒班和临时任务，其实际运输时间可以在允许范围内向后调整。
此外，8月9日需求量相对较低也可能受到规划期边界的影响，不能直接据此认为该日
所需机队规模一定较小。

综上，问题三中的不同任务在时间紧迫程度上存在显著差异。应急处置时间窗最窄、
松弛量最小，应当在调度中优先确定；增储上产次之；常规倒班和临时任务具有较大的
调整空间，可用于平衡不同日期的运输压力。上述分析为后续构建带时间窗、任务优先级
和机队资源约束的运输调度模型提供了依据。
\subsection{评价指标}
总燃油消耗按全部航段求和；座位利用率定义为
\begin{equation}
  U=\frac{\sum_r\sum_{h=0}^{m_r}L_{rh}d_{v_{rh},v_{r,h+1}}}
  {\sum_r\sum_{h=0}^{m_r}s_{k(r)}d_{v_{rh},v_{r,h+1}}}.
  \label{eq:utilization}
\end{equation}
人员在途时间从其上机点的离开时刻计至终点到达时刻，包括中途绕行、停靠和等待。

\subsection{独立仿真器与校验器}
每条路线按“到达--余油校验--下客--上客--容量校验--加油/等待--离开”递推状态。优化器只负责提出候选方案，最终指标一律由独立校验器从导出的 CSV 重新计算，从而避免算法缓存与提交结果不一致。

\begin{table}[htbp]
  \centering
  \caption{核心可行性检查清单}
  \label{tab:validator}
  \begin{tabularx}{\linewidth}{p{0.23\linewidth}X}
    \toprule
    检查类别 & 主要检查内容 \\
    \midrule
    路线结构 & 首末同机场、中途无机场、海上着陆不超过 5 次、编号连续 \\
    人员分配 & 每名必运人员恰好一次、OD 匹配、先取后送、首次到终点即下机 \\
    飞行状态 & 每航段不超座、每一到达点满足安全余油、加油点合法 \\
    Q3 排班 & 时间窗、06:00--18:00 起飞、20:00 前返场、30 min 周转、飞机所属机场 \\
    一致性 & CSV 表头与空值规则、五项指标从最终文件复算、论文数字一致 \\
    \bottomrule
  \end{tabularx}
\end{table}

\section{问题一：单向出海运输}
\label{sec:q1}

\subsection{问题结构与求解思路}

问题一要求完成1600名人员的单向出海运输，其中247名人员绑定具体机场，
其余1353名人员的起点为\texttt{LAND}，可以从任意一座机场出发。
本问不考虑人员时间窗和具体飞机数量限制，但仍需同时决定LAND人员的实际机场、
架次所用机型、海上设施访问顺序、人员分批方案和加油选择。

由于人员仅在陆地机场上机，并在其目标设施下机，任一架次的机上人数沿路线单调不增。
因此，问题一可视为一个带灵活仓库选择、异构机型、容量约束、燃油补给约束和停靠次数
约束的多仓配送问题。本文采用“单设施容量动态规划（B0）--广义节约合并（B1）--
带整数规划修复的关联度引导ALNS及标准ALNS反馈精修（B2）”的三级方法。B0保证始终具有可提交底解，B1压缩
重复往返，B2通过整条路线破坏和精确修复重新决定人员分批、LAND机场、路线、机型与
加油方案；每一级输出均由独立校验器从CSV逐航段复算。

\subsection{解的表示与路线评价}

将具有相同起点约束和目标设施的人员聚合为需求批次集合 $\mathcal{B}^{(1)}$。
对任意批次 $b\in\mathcal{B}^{(1)}$，记其人数为 $Q_b$，目标设施为 $f_b$，允许使用的机场集合为
\begin{equation}
  A_b=
  \begin{cases}
    A, & \text{若批次 $b$ 的起点为\texttt{LAND}},\\
    \{a_b\}, & \text{若批次 $b$ 绑定机场 $a_b$}.
  \end{cases}
  \label{eq:q1-compatible-airports}
\end{equation}

一条架次路线记为
\begin{equation}
  r=(v_{r0},v_{r1},\ldots,v_{r,m_r},v_{r,m_r+1}),
  \label{eq:q1-route}
\end{equation}
其中 $v_{r0}=v_{r,m_r+1}=a(r)\in A$，中间节点均属于海上设施集合 $F$，
且 $m_r\le 5$。中间节点既可以是实际送达人员的服务设施，也可以是仅用于满足燃油
约束的技术停靠点。记 $k(r)$ 为路线 $r$ 使用的机型，$\eta_{br}$ 为批次 $b$
分配到路线 $r$ 的人数。人员覆盖和机场兼容约束为
\begin{equation}
  \sum_r\eta_{br}=Q_b,
  \qquad
  \eta_{br}>0\Rightarrow
  \bigl(a(r)\in A_b,\ f_b\in r\bigr).
  \label{eq:q1-demand-cover}
\end{equation}

设 $h_r(f)$ 为路线 $r$ 第一次访问设施 $f$ 的停靠序号。由于同一人员到达其目标设施
时立即下机，路线离开第 $h$ 个海上停靠点后的载客量递推为
\begin{equation}
  L_{r0}=\sum_{b\in\mathcal{B}^{(1)}}\eta_{br},
  \qquad
  L_{rh}=L_{r,h-1}-
  \sum_{b:\,h_r(f_b)=h}\eta_{br},
  \quad h=1,\ldots,m_r.
  \label{eq:q1-load}
\end{equation}
因问题一不存在海上上客，容量约束等价于
\begin{equation}
  0\le L_{rh}\le s_{k(r)},
  \qquad h=0,1,\ldots,m_r.
  \label{eq:q1-capacity}
\end{equation}
其中初始载客量 $L_{r0}$ 是路线上的最大载客量。

记 $\xi_{rh}\in\{0,1\}$ 为第 $h$ 个海上停靠点的加油决策。
飞机从机场满油起飞，即
$G_{r0}^{\mathrm{dep}}=G_{k(r)}^{\max}$。逐航段燃油状态满足
\begin{equation}
  G_{r,h+1}^{\mathrm{arr}}
  =G_{rh}^{\mathrm{dep}}
  -\beta_{k(r)}d_{v_{rh},v_{r,h+1}}
  \ge G_{k(r)}^{\min},
  \qquad h=0,\ldots,m_r,
  \label{eq:q1-fuel-arrival}
\end{equation}
且对于海上停靠点 $h=1,\ldots,m_r$，
\begin{equation}
  G_{rh}^{\mathrm{dep}}=
  \begin{cases}
    G_{k(r)}^{\max}, & v_{rh}\in R,\ \xi_{rh}=1,\\
    G_{rh}^{\mathrm{arr}}, & \xi_{rh}=0.
  \end{cases}
  \label{eq:q1-fuel-departure}
\end{equation}
若 $v_{rh}\notin R$，则必须有 $\xi_{rh}=0$。问题一无需主动等待，故海上停靠时间为
\begin{equation}
  \tau_{rh}=10+10\xi_{rh},
  \qquad h=1,\ldots,m_r.
  \label{eq:q1-dwell}
\end{equation}

据此，路线 $r$ 的飞机使用时间和燃油消耗量分别为
\begin{align}
  T_r^{\mathrm{air}}
  &=\sum_{h=0}^{m_r}t_{v_{rh},v_{r,h+1}}^{k(r)}
    +\sum_{h=1}^{m_r}\tau_{rh},
  \label{eq:q1-route-time}\\
  F_r
  &=\sum_{h=0}^{m_r}\beta_{k(r)}d_{v_{rh},v_{r,h+1}}.
  \label{eq:q1-route-fuel}
\end{align}
批次 $b$ 中人员在路线 $r$ 上的单人在途时间为
\begin{equation}
  T_{br}^{\mathrm{pass}}
  =\sum_{h=0}^{h_r(f_b)-1}t_{v_{rh},v_{r,h+1}}^{k(r)}
  +\sum_{h=1}^{h_r(f_b)-1}\tau_{rh},
  \label{eq:q1-passenger-time}
\end{equation}
其中不计入人员到达目标设施后的停靠时间。

对一个完整方案 $S$，定义统一最小化方向下的词典序目标向量
\begin{equation}
  Z(S)=\left(
  T_{\mathrm{air}},\
  T_{\mathrm{passenger}},\
  N_{\mathrm{flight}},\
  F_{\mathrm{fuel}},\
  -U
  \right),
  \label{eq:q1-objective}
\end{equation}
其中各指标均由最终路线重新计算。算法首先严格最小化总飞机使用时间；只有当前一级
指标相等时，才依次比较人员总在途时间、架次数、燃油消耗量和座位利用率。

\subsection{覆盖全部需求的可行初解}

\subsubsection{LAND人员的初始机场分配}

对于绑定具体机场的需求，保持其机场不变。对于起点为\texttt{LAND}、目标为设施 $f$
的需求，首先对每个机场 $a\in A$ 和机型 $k\in K$，利用燃油标签搜索计算一条
“机场 $a$ 出发--访问设施 $f$--返回机场 $a$”的最短燃油可行闭合路线。令
$\omega_{af}$ 为该OD分配至机场 $a$ 的人数，则外层枚举满足
\begin{equation}
  \sum_{a\in A}\omega_{af}=Q_f^{\mathrm{LAND}},\qquad
  \omega_{af}\in\mathbb Z_{\ge 0},
  \label{eq:q1-initial-airport}
\end{equation}
并以各机场容量动态规划目标向量之和的词典序最小者作为初始分流方案。这样允许同一设施
的LAND人员拆分到不同机场，避免“全分给最近机场”造成余量架次；B2修复阶段仍可再次
跨机场重分配LAND需求，而指定机场人员始终不可改变机场。

\subsubsection{单设施安全底解}

对于每个已确定机场的聚合OD，分别枚举三种机型。若某机型不能直接闭合往返，
则在最多5次海上着陆范围内插入技术停靠点和必要加油点。由于单个OD人数不超过较小规模，
本文使用容量动态规划确定不同机型架次的组合，使所有人员得到覆盖。令
$\lambda$ 表示尚有 $n$ 人时新增架次的实际装载人数，
$\Delta Z_{af}^{k}(\lambda)$ 表示该架次产生的词典序指标增量，则
\begin{equation}
  \Phi_{af}(n)=
  \min_{\substack{k\in K,\;1\le\lambda\le\min\{n,s_k\}}}^{\mathrm{lex}}
  \left\{
    \Phi_{af}(n-\lambda)+\Delta Z_{af}^{k}(\lambda)
  \right\},
  \qquad \Phi_{af}(0)=\boldsymbol{0},
  \label{eq:q1-batch-dp}
\end{equation}
其中 $n$ 为尚未运输的人数。枚举 $\lambda$ 而不是预设为满载，可以在三种机型容量间
寻找更优整数拆分，统一处理满载与余量架次，避免按某一种机型座位数机械拆批。由此得到
覆盖全部1600人的单设施安全底解，并立即交由独立校验器验证。

\subsubsection{广义节约合并}

在安全底解基础上，使用Clarke--Wright节约思想\cite{clarke1964savings}合并同一机场的
不同路线。由于本题包含异构机型和燃油补给，经典距离节约值不足以评价真实改进，
本文定义两条路线 $r_1,r_2$ 的广义节约值为
\begin{equation}
  \operatorname{Save}(r_1,r_2)
  =T_{r_1}^{\mathrm{air}}+T_{r_2}^{\mathrm{air}}
  -\min_{r\in\mathcal{M}(r_1,r_2)}T_r^{\mathrm{air}},
  \label{eq:q1-generalized-saving}
\end{equation}
其中 $\mathcal{M}(r_1,r_2)$ 表示保持机场相同、覆盖原有人员且满足容量、燃油和
5次海上着陆约束的合并路线集合。计算合并路线时同时枚举三种机型、服务设施顺序和
加油方案，不采用“先合并路线、再被动选择机型”的顺序。算法按广义节约值从大到小
尝试合并；当飞机时间节约值为0时，仅接受能够改善后续词典序指标的合并。
直至不存在使完整目标向量改进的可行合并，得到ALNS的初始方案
$S_{\mathrm{init}}^{(1)}$。

\subsection{机型与加油联合评价的标签算法}

每次路线插入、合并或重排后，都需要重新判断其燃油可行性。对于给定机场、机型和
海上停靠序列，本文采用标签算法决定各加油点是否加油。第 $h$ 个停靠位置的一个标签
记为
\begin{equation}
  \ell_h=(v_h,h,m_h,T_h,G_h,F_h),
\end{equation}
其中 $v_h$ 为当前位置，$m_h$ 为已使用的海上停靠次数，$T_h$、$G_h$ 和 $F_h$ 分别为
累计时间、离开该点的剩余燃油和累计耗油量。初始标签对应机场、0次停靠、0 min、
$G_k^{\max}$ 剩余燃油和0耗油。对下一航段，若到达油量低于
$G_k^{\min}$，则删除该标签；若下一节点
不是加油设施，仅生成“不加油”标签；若下一节点属于 $R$，则分别生成“不加油”和
“加满油”两个标签。

在位置与停靠次数相同的状态下，若两个标签满足
\begin{equation}
  T_h^{(1)}\le T_h^{(2)},
  \qquad
  G_h^{(1)}\ge G_h^{(2)},
  \qquad
  F_h^{(1)}\le F_h^{(2)},
  \label{eq:q1-label-dominance}
\end{equation}
则标签1支配标签2，后者可以删除。最终仅保留能够安全返回原机场的标签，并选择
使用时间最小者；若使用时间相同，再比较燃油状态和后续目标增量。该子程序避免了
“经过加油点就必定加油”的保守规则，同时保证每次路线变化后的机型与加油决策同步更新。

\subsection{基于整数规划修复的自适应大邻域搜索}

广义节约算法只能合并总人数不超过单机容量的路线，无法改变已形成的人员分批，也难以
修正早期LAND机场分配。为此，本文以 $S_{\mathrm{init}}^{(1)}$ 为初始解，构造带整数规划
修复的ALNS\citep{shaw1998lns,ropke2006alns,pisinger2010lns}。算法保留前述精确路线
评价与燃油标签搜索，不以欧氏距离或近似时长替代最终航段评价。

\subsubsection{五类整路线破坏算子}

每次迭代移除若干完整路线，将其中乘客恢复为按“起点约束--目的设施”聚合的待修复需求组。
程序实际实现以下五种算子：
\begin{enumerate}[label=\arabic*.]
  \item \textbf{随机路线移除：}保持搜索多样性；
  \item \textbf{相关路线移除：}联合移除机场相同、服务设施相近的路线；
  \item \textbf{低利用率路线移除：}优先重组座位公里利用率低的余量路线；
  \item \textbf{单位人员时间最差路线移除：}优先处理飞机时间与载客人数比值大的路线；
  \item \textbf{LAND重指派移除：}优先选择LAND人员较多的路线，使修复阶段可以换机场。
\end{enumerate}

\subsubsection{候选路线生成与精确修复}

设本轮待修复需求组集合为 $\mathcal G_D$，候选路线集合为
$\mathcal R_D^{\mathrm{cand}}$。每条候选路线 $c$ 同时确定机场 $a(c)$、机型 $k(c)$
和至多两个实际服务设施的访问顺序；标签搜索可以继续插入技术经停与加油点，但总海上
着陆数不得超过5。令 $Q_g^D$ 为需求组人数，$\zeta_c$ 为路线 $c$ 的使用架次数，
$\chi_{gc}$ 为组 $g$ 分配给路线 $c$ 的人数，则修复模型满足
\begin{align}
  \sum_{c\in\mathcal R_D^{\mathrm{cand}}}\chi_{gc}&=Q_g^D,
  &&g\in\mathcal G_D,\label{eq:q1-repair-cover}\\
  \sum_{g\in\mathcal G_D}\chi_{gc}&\le s_{k(c)}\zeta_c,
  &&c\in\mathcal R_D^{\mathrm{cand}}.\label{eq:q1-repair-capacity}
\end{align}
若目的设施不在路线中，或指定机场与 $a(c)$ 不一致，则固定 $\chi_{gc}=0$；LAND需求
允许分配至任意机场。为防止候选路线包含无人服务的冗余停靠，对路线中的每个服务设施
$f$ 还要求
\begin{equation}
  \sum_{g:\,f_g=f}\chi_{gc}\ge\zeta_c.
  \label{eq:q1-repair-service}
\end{equation}
修复模型首先最小化
\begin{equation}
  \min\sum_{c\in\mathcal R_D^{\mathrm{cand}}}T_c^{\mathrm{air}}\zeta_c,
  \label{eq:q1-repair-objective}
\end{equation}
固定一级最优值后，依次最小化
$\sum_{g,c}\vartheta_{cg}\chi_{gc}$、$\sum_c\zeta_c$ 和
$\sum_c F_c\zeta_c$，分别对应人员在途时间、架次数和燃油消耗。聚合解按人员编号回填，
并经内存路线评价器和最终独立Validator双重检查。

\subsubsection{自适应权重与接受准则}

设算子 $\ell$ 在第 $\gamma$ 个统计周期内被调用
$n_{\ell}^{(\gamma)}$ 次，累计奖励分数为 $\sigma_{\ell}^{(\gamma)}$，
则其权重更新为
\begin{equation}
  w_{\ell}^{(\gamma+1)}
  =(1-\rho)w_{\ell}^{(\gamma)}
  +\rho\frac{\sigma_{\ell}^{(\gamma)}}{\max\{1,n_{\ell}^{(\gamma)}\}},
  \label{eq:q1-alns-weight}
\end{equation}
其中 $\rho\in(0,1)$ 为反应系数。产生全局最优解、改进当前解和被接受的劣解分别获得
不同分值，使有效算子在后续迭代中具有更高的选择概率。

若候选解 $S'$ 在词典序意义下优于当前解 $S$，则直接接受。若候选解较差，令
$j^\star$ 为目标向量中第一个不同的分量，并定义相对劣化量
\begin{equation}
  \Delta(S',S)=
  \frac{Z_{j^\star}(S')-Z_{j^\star}(S)}
  {\max\{1,|Z_{j^\star}(S)|\}},
  \label{eq:q1-deterioration}
\end{equation}
则以概率
\begin{equation}
  P_{\mathrm{acc}}
  =\exp\left(-\frac{\max\{0,\Delta(S',S)\}}{\Theta}\right)
  \label{eq:q1-acceptance}
\end{equation}
接受候选解，其中温度按照 $\Theta\leftarrow\alpha\Theta$ 逐步下降。
该策略允许算法暂时越过局部最优，但历史最优解始终按式~\eqref{eq:q1-objective}
的词典序规则保存。

\subsubsection{关联度引导与标准ALNS反馈精修}

公平比较采用相同300 s预算和5个随机种子。Combined R3的最好值、中位数和均值分别为
15025、15149和15142.0 min，而标准ALNS V1分别为15118、15185和15208.4 min；R3在
同种子比较中4/5胜出，并使筛选阶段精确评价调用量平均减少57.0\%。因此本文采用R3负责
搜索新盆地，但不把相关度评分固化为最终求解器。R3扩展搜索先把15025 min降至14772 min，
随后以该解热启动纯标准ALNS A2，在82.207 s处取得14770 min。该解不再直接作为终点，
而是连同不同种子、不同引导方式和局部强化产生的精英路线一起送入路线池主问题。

\subsubsection{精英路线池与精确主问题重组}

令$\mathcal R^{\mathrm{pool}}$为跨轨迹收集并按“机场--机型--停靠序列--加油位置”去重后的
物理路线集合。由于同一物理路线可以承载不同人员组合，程序进一步把路线与可行人员分配
绑定为allocated route pattern $p\in\mathcal P$。令$z_p\in\{0,1\}$表示是否采用模式$p$，
$a_{bp}$表示模式$p$覆盖批次$b$的人数，则精确主问题为
\begin{align}
  \min\quad &\sum_{p\in\mathcal P}T_p^{\mathrm{air}}z_p,
  \label{eq:q1-pool-master-objective}\\
  \text{s.t.}\quad &\sum_{p\in\mathcal P}a_{bp}z_p=Q_b,
  &&b\in\mathcal B^{(1)},\label{eq:q1-pool-master-cover}\\
  &z_p\in\{0,1\},&&p\in\mathcal P.
\end{align}
路线模式在入池前已经通过容量、燃油、机场绑定与停靠次数检查，故主问题只负责跨轨迹重组
和人员完整覆盖。第四轮路线池含95个来源、391条不同物理路线和996个分配模式；整数主问题
把14770 min降至14730 min。再以该解热启动标准ALNS，一级目标保持14730 min，同时把人员
在途时间由122494 min降至121363 min，形成最终词典序胜者。另得到14732 min、88架次的
可行解，但因首要目标多2 min，严格按词典序不予替代。

\subsubsection{多阶段参数与主流程}

第一阶段侧重快速清除明显低利用率路线，第二阶段扩大破坏规模并强化重组。实际参数如
表~\ref{tab:q1-alns-parameters}所示。单条候选路线最多含两个实际服务设施，但标签算法
可插入技术经停点，总海上停靠次数仍不超过5。

\begin{table}[htbp]
  \centering
  \caption{问题一融合搜索的关键阶段}
  \label{tab:q1-alns-parameters}
  \begin{tabular}{lrrrrr}
    \toprule
    阶段 & 种子 & 预算 & 引导方式 & 最好值/min & 作用\\
    \midrule
    标准ALNS公平对照 & 0--4 & $5\times300$ s & 标准相关移除 & 15118 & 对照\\
    Combined R3公平搜索 & 0--4 & $5\times300$ s & 距离+Context V2 & 15025 & 选入\\
    R3扩展搜索 & 3 & 808.779 s & 距离+Context V2 & 14772 & 跨盆地\\
    标准A2热启动精修 & 4 & 604.658 s & 标准ALNS & 14770 & 路线池控制解\\
    第四轮精确主问题 & -- & 有限时限 & 路线池MILP & \textbf{14730} & 一级最优候选\\
    标准ALNS教育 & 21 & 有限时限 & 热启动精修 & \textbf{14730} & 最终词典序解\\
    \bottomrule
  \end{tabular}
\end{table}

问题一的完整优化流程如下：
\begin{enumerate}[label=\textbf{步骤\arabic*.},leftmargin=4.5em]
  \item 读取预处理后的Q1需求、航段特征和燃油可达性数据；
  \item 固定指定机场需求，并按式~\eqref{eq:q1-initial-airport}为LAND需求选择初始机场；
  \item 通过容量动态规划和燃油标签算法生成覆盖全部人员的单设施安全底解；
  \item 按式~\eqref{eq:q1-generalized-saving}执行广义节约合并，得到初始方案
  $S_{\mathrm{init}}^{(1)}$；
  \item 按当前权重选择一种整路线破坏算子，形成待修复需求组；
  \item 枚举机场、机型和服务顺序，调用标签搜索补足技术经停与加油方案；
  \item 求解式~\eqref{eq:q1-repair-cover}--\eqref{eq:q1-repair-objective}的整数修复模型；
  \item 由内存路线评价器检查人员覆盖、机场匹配、容量、油量、停靠和闭合约束；
  \item 按模拟退火准则决定是否接受候选解，并更新历史最优解和算子得分；
  \item 达到最大运行时间、最大迭代次数或连续若干轮无改进时停止；
  \item 在同预算多种子试验中比较标准ALNS与Combined R3，保留严格词典序更优者；
  \item 延长R3搜索以突破原有盆地，再将其历史最优解反馈给标准A2作热启动精修；
  \item 汇集多轨迹精英路线并生成带人员分配的路线模式，求解式~\eqref{eq:q1-pool-master-objective}--\eqref{eq:q1-pool-master-cover}；
  \item 将主问题最优候选反馈给标准ALNS作education，在一级目标不变时改善后续指标；
  \item 将最终聚合解按人员编号回填，导出两个CSV并用独立Validator复算。
\end{enumerate}

各阶段耗时分别记录，不把公平试验、路线池构建、主问题与最终education混写成一次运行时间。
所有比较均使用独立Validator通过的完整CSV，未通过候选不进入统计。

\subsection{算法实现与输出文件}

问题一程序按支撑材料实际名称组织：\path{src/solver/baseline.py}实现容量动态规划与LAND分流，
\path{src/solver/improve.py}实现广义节约，\path{src/solver/alns.py}实现五类破坏算子、
候选路线生成、SciPy MILP修复、自适应权重和模拟退火接受；
\path{src/solver/relatedness.py}实现原始距离、冻结共识读取和Context V2候选排序；
\path{src/solver/q1\_or.py}实现精英路线池、allocated-pattern主问题、LP放松、定向邻域和路线审计；
\path{scripts/05\_run\_alns\_multiseed.py}实现公平多种子试验，
\path{scripts/06\_finalize\_q1\_relatedness.py}按严格词典序固化最终方案；
\path{scripts/07\_run\_q1\_route\_pool\_master.py}--\path{scripts/13\_finalize\_q1\_or.py}
依次执行路线池主问题、定向强化、精英重组、路径重连、修复对照、学习数据审计和最终固化；
\path{src/solver/technical\_stops.py}负责技术经停与加油标签搜索，
\path{src/solver/evaluator.py}负责内存路线评价，
\path{src/solver/exporter.py}负责结果导出，
\path{src/validation/solution\_validator.py}从最终CSV独立复算。
基础入口为\path{scripts/04\_solve\_q1\_alns.py}，中文阶段图由
\path{scripts/plot\_integrated\_improvements.py}从比较表生成。程序实际保存
\texttt{q1-routes.csv}、\texttt{q1-assignments.csv}、\texttt{metrics.json}、
\texttt{validator.json}、\texttt{winning\_config.json}和\texttt{method\_metadata.json}。

\subsection{结果展示与算法分析}

表~\ref{tab:q1-results}比较从安全底解到融合最终解的主要阶段。各阶段均覆盖1600人并通过独立校验，
所以改进来自路线组织、人员重新分批、LAND机场分流、机型与加油决策，而非放松约束。

\begin{table}[htbp]
  \centering
  \caption{问题一不同算法的结果比较}
  \label{tab:q1-results}
  \resizebox{\linewidth}{!}{%
  \begin{tabular}{lrrrrrr}
    \toprule
    方法 & 飞机时间/min & 人员在途/min & 架次数 & 燃油/kg & 利用率/\% & 运行时间/s \\
    \midrule
    单设施容量动态规划（B0） & 17222 & 120609 & 110 & 142601.0 & 43.047 & -- \\
    广义节约（B1）           & 15743 & 120396 & 100 & 128562.3 & 46.757 & -- \\
    经典VND                    & 15371 & 120870 & 95 & 123081.7 & 48.226 & 149.000 \\
    标准ALNS公平对照           & 15118 & 122061 & 92 & 121682.9 & 48.446 & 329.256 \\
    Combined R3公平最优        & 15025 & 121524 & 91 & 121382.4 & 49.056 & 310.083 \\
    R3扩展搜索                 & 14772 & 121652 & 89 & 118645.9 & 49.346 & 808.779 \\
    标准A2路线池控制解         & 14770 & 120845 & 89 & 118640.1 & 49.339 & 604.658 \\
    路线池精确主问题           & \textbf{14730} & 122494 & 89 & \textbf{118624.4} & 49.404 & -- \\
    主问题+标准ALNS教育        & \textbf{14730} & \textbf{121363} & \textbf{89} &
                                  \textbf{118624.4} & \textbf{49.430} & -- \\
    \bottomrule
  \end{tabular}}
\end{table}

本文严格优先最小化总飞机使用时间，因此一级目标下降时允许人员在途时间略增；这正是
式~\eqref{eq:q1-objective}的词典序含义。融合算法相对B1的一级目标改进率为
\begin{equation}
  I_{\mathrm{Q1}}
  =\frac{15743-14730}{15743}\times100\%=6.435\%.
  \label{eq:q1-improvement}
\end{equation}
同时减少11个架次和9937.9 kg燃油，座位利用率提高2.673个百分点。

图~\ref{fig:q1-convergence}由各阶段已验证比较表直接绘制，展示容量底解、广义节约、
经典VND、标准ALNS、关联度引导和最终反馈精修的完整改进链。

\begin{figure}[htbp]
  \centering
  \IfFileExists{figures/q1_integrated_improvement.pdf}
    {\includegraphics[width=0.88\linewidth]{figures/q1_integrated_improvement.pdf}}
    {\figplaceholder{Q1 ALNS收敛曲线：当前解与历史最优解}}
  \caption{问题一精英路线池主问题与标准ALNS教育的阶段改进}
  \label{fig:q1-convergence}
\end{figure}

最终89条路线中，T1、T2、T3分别占0、26、63架次；A01、A02、A03分别承担45、20、24
架次。54条路线服务1个实际设施，35条服务2个实际设施；标签算法共安排32次加油。
每架载客数为15--19人，平均17.978人，说明融合算法的主要收益是消除低装载余量路线并将其
重新组合，而不是生成复杂的长停靠路线。

为判断解的绝对质量，表~\ref{tab:q1-lower-bounds}将最终结果与各指标的理论界比较。
这些界分别来自放松容量共享、返场、燃油或路线耦合后的独立问题，通常不能由同一方案
同时达到，故只能作为保守参照。

\begin{table}[htbp]
  \centering
  \caption{问题一最终结果与理论界对照}
  \label{tab:q1-lower-bounds}
  \begin{tabular}{lrrr}
    \toprule
    指标 & 理论界 & B2结果 & 相对间隙/\% \\
    \midrule
    总飞机使用时间/min & 13337 & 14730 & 10.445 \\
    人员总在途时间/min & 95328 & 121363 & 27.311 \\
    总架次数 & 85 & 89 & 4.706 \\
    总燃油消耗/kg & 113878.73 & 118624.4 & 4.167 \\
    座位利用率/\% & 100（上界） & 49.430 & 50.570个百分点 \\
    \bottomrule
  \end{tabular}
\end{table}

其中一级目标的保守最优性间隙为
\begin{equation}
  \operatorname{Gap}_{\mathrm{Q1}}
  =\frac{14730-13337}{13337}\times100\%=10.445\%.
  \label{eq:q1-gap}
\end{equation}
该值不是算法真实误差的精确值，而是由原问题放松下界给出的保守区间；下界越松，区间越大。
此外，第四轮路线池LP值为14208.637 min，与最终整数值的池内参考间隙为
$ (14730-14208.637)/14208.637=3.669\%$。该LP只放松当前996个分配模式的整数性，
既没有覆盖所有潜在路线，也不能替代13337 min的原问题全局有效下界。

表~\ref{tab:q1-operator-stats}汇总14770路线池控制解对应的标准A2热启动日志。相关路线移除
产生3次历史最优，是该阶段唯一直接刷新最好解的算子；其余算子维持可行盆地与邻域多样性。
最终14730的额外40 min收益来自跨轨迹路线池重组，故该日志不能单独代表完整算法贡献。

\begin{table}[htbp]
  \centering
  \caption{问题一14770控制解的标准ALNS算子统计}
  \label{tab:q1-operator-stats}
  \begin{tabular}{lrrrr}
    \toprule
    破坏算子 & 调用次数 & 接受次数 & 改进当前解 & 产生历史最优 \\
    \midrule
    LAND重指派 & 18 & 16 & 0 & 0 \\
    低利用率路线 & 19 & 19 & 0 & 0 \\
    随机路线 & 23 & 23 & 0 & 0 \\
    相关路线 & 31 & 29 & 4 & 3 \\
    单位人员时间最差路线 & 17 & 17 & 0 & 0 \\
    \bottomrule
  \end{tabular}
\end{table}

最终\texttt{validator.json}给出\texttt{valid=true}、\texttt{issue\_count=0}，且内存指标
与从两个最终CSV独立复算的指标完全一致。全量89条路线和1600条人员分配记录放入支撑
材料，正文仅保留汇总指标、阶段改进和理论界对照。

\section{问题二：出海、海返与穿梭联合运输}
\label{sec:q2}

\subsection{问题结构与总体思路}

问题二包含4000条人员运输需求，其中出海、海返和设施间穿梭需求分别为
1600、1600和800条。与问题一相同，本问不考虑时间窗和具体飞机数量限制，
但每名人员必须由同一架次完成运输，不得换乘；每个架次仍须从某座机场起飞并返回
同一机场，且海上着陆次数不超过5次。

问题二的主要优化空间来自座位在同一架次内的连续复用：飞机在某设施送下出海人员后，
刚释放的座位可以立即接运海返或穿梭人员；穿梭人员到达目标设施下机后，座位又可继续
服务后续需求。因此，一架次累计运输人数可以大于机型座位数，真正需要约束的是每个航段
离开停靠点后的同时在机人数。本文将问题二建模为带灵活机场选择、异构机型、燃油补给和
最多5次海上着陆约束的多仓成对取送车辆路径问题。本文复用Q1燃油标签与路线评价器，
采用“三类分开运输基线--双设施候选路线整数主问题--自适应三设施列扩展”的三级方法。
与逐批次Regret插入相比，该方法直接在聚合OD层同时决定路线使用次数和人员流量，并用
逐航段容量约束刻画同一座位在一次架次中的连续复用。

\subsection{联合取送路线模型}

\subsubsection{需求批次与机场兼容性}

将起点和终点相同的人员聚合为需求批次集合 $\mathcal{B}^{(2)}$。
对任意批次 $b\in\mathcal{B}^{(2)}$，
记其人数为 $Q_b$，起点和终点分别为 $o_b$ 和 $d_b$。允许该批次使用的机场集合定义为
\begin{equation}
  A_b=
  \begin{cases}
    \{a\}, & o_b=a\in A\ \text{或}\ d_b=a\in A,\\
    A, & \text{其余情况}.
  \end{cases}
  \label{eq:q2-compatible-airports}
\end{equation}
其中，当陆地端为\texttt{LAND}时可以选择任意机场；设施间穿梭需求同样可以由任意机场
闭合的架次承担。最终方案生成时，再将聚合数量还原为唯一人员编号。

问题二的一条架次仍记为
\begin{equation}
  r=(v_{r0},v_{r1},\ldots,v_{r,m_r},v_{r,m_r+1}),
  \qquad v_{r0}=v_{r,m_r+1}=a(r),\quad m_r\le5.
  \label{eq:q2-route}
\end{equation}
设 $\eta_{br}$ 为批次 $b$ 分配到路线 $r$ 的人数，则人员覆盖约束为
\begin{equation}
  \sum_r\eta_{br}=Q_b,
  \qquad b\in\mathcal{B}^{(2)},
  \label{eq:q2-demand-cover}
\end{equation}
且 $\eta_{br}>0$ 时必须有 $a(r)\in A_b$。

\subsubsection{取送位置与禁止换乘约束}

记 $\pi_{br}$ 和 $\delta_{br}$ 分别为批次 $b$ 在路线 $r$ 上的上机和下机序号。
对出海、海返和穿梭需求，分别按照下式确定：
\begin{equation}
  \pi_{br}=
  \begin{cases}
    0, & o_b\in A\cup\{\mathrm{LAND}\},\\
    \min\{h:v_{rh}=o_b\}, & o_b\in F,
  \end{cases}
  \label{eq:q2-pickup-order}
\end{equation}
\begin{equation}
  \delta_{br}=
  \begin{cases}
    m_r+1, & d_b\in A\cup\{\mathrm{LAND}\},\\
    \min\{h>\pi_{br}:v_{rh}=d_b\}, & d_b\in F.
  \end{cases}
  \label{eq:q2-delivery-order}
\end{equation}
因此，对所有 $\eta_{br}>0$ 的批次均有
\begin{equation}
  \pi_{br}<\delta_{br}.
  \label{eq:q2-precedence}
\end{equation}
式~\eqref{eq:q2-delivery-order}同时保证人员在上机后第一次到达其终点时立即下机。
同一批次的上机与下机位置属于同一条路线，故禁止换乘约束自动得到满足。

\subsubsection{逐站载荷递推}

记第 $h$ 个停靠点的下机人数和上机人数分别为
\begin{equation}
  N_{rh}^{\mathrm{off}}=\sum_{b:\,\delta_{br}=h}\eta_{br},
  \qquad
  N_{rh}^{\mathrm{on}}=\sum_{b:\,\pi_{br}=h}\eta_{br}.
  \label{eq:q2-pickup-delivery-count}
\end{equation}
机场起飞时的初始载荷为 $L_{r0}=N_{r0}^{\mathrm{on}}$。
按照题目规定的“先下后上”顺序，
路线离开第 $h$ 个海上停靠点后的载荷为
\begin{equation}
  L_{rh}=L_{r,h-1}-N_{rh}^{\mathrm{off}}+N_{rh}^{\mathrm{on}},
  \qquad h=1,\ldots,m_r.
  \label{eq:q2-load}
\end{equation}
容量约束为
\begin{equation}
  0\le L_{rh}\le s_{k(r)},
  \qquad h=0,1,\ldots,m_r.
  \label{eq:q2-capacity}
\end{equation}
返场前还须满足 $L_{r,m_r}=N_{r,m_r+1}^{\mathrm{off}}$，从而保证所有海返人员在机场下机后
机上人数归零。该递推说明不能用“架次累计运输人数不超过座位数”代替逐航段容量检查。

\subsubsection{路线时间、燃油与人员在途时间}

路线的飞行时间、停靠时间和燃油状态沿用问题一的定义。每次到达均须满足安全余油，
仅可在指定的8座设施加满油；不加油停靠至少10 min，加油停靠至少20 min。
对于分配到路线 $r$ 的批次 $b$，其单人在途时间为
\begin{equation}
  T_{br}^{\mathrm{pass}}
  =\sum_{h=\pi_{br}}^{\delta_{br}-1}
  t_{v_{rh},v_{r,h+1}}^{k(r)}
  +\sum_{h=\pi_{br}+1}^{\delta_{br}-1}\tau_{rh}.
  \label{eq:q2-passenger-time}
\end{equation}
该式不计人员上机点离开前的停靠时间，也不计其到达目标设施后的停靠时间。
问题二仍按照
\begin{equation}
  \min_{\mathrm{lex}}
  \left(T_{\mathrm{air}},T_{\mathrm{passenger}},
  N_{\mathrm{flight}},F_{\mathrm{fuel}},-U\right)
  \label{eq:q2-objective}
\end{equation}
进行词典序优化，且所有指标由最终CSV重新仿真计算。

\subsection{候选闭合路线生成}

Q2的1600条出海需求与Q1完全一致，因此首先读取Q1最优路线，将其中实际服务设施顺序
加入候选池；但Q2主问题不固定这些路线，允许海返和穿梭客流改变机场、机型和路线组合。
为避免枚举 $52P_5$ 个停靠序列，本文分四类生成服务序列：
\begin{enumerate}[label=\arabic*.]
  \item 52条单设施序列，用于保证每个设施均有安全底解；
  \item Q2中真实存在的设施间穿梭OD序列 $(o_b,d_b)$；
  \item 每个设施的3个空间近邻及10个高需求设施的有序配对；
  \item Q1最终路线中的实际服务顺序及其逆序。
\end{enumerate}
由此得到357条服务序列。对每条序列分别枚举3座机场和3种机型，并调用问题一的燃油
标签算法插入必要技术经停与加油点；删除燃油不可行或超过5次海上着陆的路线后，得到
3116条候选闭合路线集合 $\mathcal C$。

对候选路线 $c$ 和批次 $b$，程序按式~\eqref{eq:q2-pickup-order}--
\eqref{eq:q2-delivery-order}确定取送位置。若指定机场不匹配、节点不存在或取送逆序，
则禁止该分配；否则计算单人在途时间 $\theta_{bc}$，并定义
\begin{equation}
  \kappa_{bch}=\mathbb I\{\pi_{bc}\le h<\delta_{bc}\},
  \label{eq:q2-onboard-indicator}
\end{equation}
表示批次 $b$ 是否占用候选路线 $c$ 的第 $h$ 个航段。

\subsection{带逐航段容量的候选路线整数主问题}

令 $y_c$ 为候选路线 $c$ 的使用架次数，$x_{bc}$ 为批次 $b$ 分配给路线 $c$ 的人数。
需求覆盖约束为
\begin{equation}
  \sum_{c\in\mathcal C}x_{bc}=Q_b,
  \qquad b\in\mathcal B^{(2)},
  \label{eq:q2-master-cover}
\end{equation}
逐航段容量约束为
\begin{equation}
  \sum_{b\in\mathcal B^{(2)}}\kappa_{bch}x_{bc}
  \le s_{k(c)}y_c,
  \qquad c\in\mathcal C,\ h=0,\ldots,m_c,
  \label{eq:q2-master-capacity}
\end{equation}
其中 $x_{bc},y_c\in\mathbb Z_{\ge0}$。式~\eqref{eq:q2-master-capacity}允许人员在中途
下机后由其他人员继续使用同一座位，因而路线累计运输人数可以超过机型座位数，但任一
航段的同时在机人数始终不超座。

一级主问题为
\begin{equation}
  \min\sum_{c\in\mathcal C}T_c^{\mathrm{air}}y_c.
  \label{eq:q2-master-primary}
\end{equation}
受固定求解时限影响，一级主问题返回当前最好可行值 $\widehat Z_1$、对偶下界
$Z_1^{\mathrm{LB}}$和相对间隙。本文不把未证明最优的$\widehat Z_1$称为全局最优值，
而是在约束
\begin{equation}
  \sum_{c\in\mathcal C}T_c^{\mathrm{air}}y_c=\widehat Z_1
  \label{eq:q2-fix-primary}
\end{equation}
下继续优化后三项目标。记燃油以$0.1$ kg为整数单位，令
\begin{equation}
  \overline N=\left\lfloor
  \frac{\widehat Z_1}{\min_{c\in\mathcal C}T_c^{\mathrm{air}}}
  \right\rfloor,
  \qquad
  \overline F=\left\lceil
  \widehat Z_1\max_{c\in\mathcal C}
  \frac{10F_c}{T_c^{\mathrm{air}}}
  \right\rceil .
  \label{eq:q2-secondary-bounds}
\end{equation}
在式~\eqref{eq:q2-fix-primary}下，任一可行解的架次数和燃油整数值分别不超过
$\overline N$和$\overline F$。据此设置
\begin{equation}
  W_F=1,\qquad
  W_N=\overline F+1,\qquad
  W_P=\overline N W_N+\overline F+1,
  \label{eq:q2-dominance-weights}
\end{equation}
并求解
\begin{equation}
  \min\quad
  W_P\sum_{b,c}\theta_{bc}x_{bc}
  +W_N\sum_cy_c
  +W_F\sum_c(10F_c)y_c.
  \label{eq:q2-secondary-objective}
\end{equation}
由于任一低优先级目标在可行域内的最大变化均小于高一级目标一个单位对应的权重，
式~\eqref{eq:q2-secondary-objective}严格保持“人员在途时间--架次数--燃油”的比较顺序，
不存在经验权重相互补偿。若第二阶段在时限内未证明最优，则保留其最好可行解；座位利用率
由最终逐航段载荷复算，只在前四项完全相同时作为方案间的最终判别量。

主问题得到的是各OD在路线类型间的聚合分配量。程序把每名乘客看作停靠序列上的一个
区间 $[\pi_{bc},\delta_{bc})$，采用区间图最优着色把乘客分配到具体座位轨道，再将每
$s_{k(c)}$ 条轨道组成一个实际架次。由于式~\eqref{eq:q2-master-capacity}限制了最大重叠
区间数，该分解一定不超过 $y_c$ 个架次，从而能够无损还原为4000条人员记录。

\subsection{精确局部MILP与多重启强化}

完整候选主问题给出19736 min可行基准后，本文不再反复求解全部264个OD的大模型，而在
当前解上执行有界破坏--精确修复。每轮由高成本、低利用率、共享设施流和LAND占比四类
算子选择相关路线集合$D$，将其人员完整取出；未被破坏的路线固定。程序围绕$D$生成
1--5个实际服务设施的有界序列，保留原路线列，并按闭合几何成本和OD覆盖潜力排序。

设局部候选路线集合为$\mathcal C_D$，局部需求为$\mathcal B_D$。修复模型仍使用
式~\eqref{eq:q2-master-cover}--\eqref{eq:q2-master-capacity}的逐航段容量结构，但只对
$D$涉及的人员和候选列建模。接受准则首先要求新局部路线完整覆盖$\mathcal B_D$且通过
精确航段评价，再要求替换后的全局词典序向量优于历史最优；模拟退火只允许接受等一级
目标的结构变化，不允许用次级目标补偿飞机时间恶化。

为降低单一初始盆地影响，程序维护质量约束精英池，并依次执行多种子延长搜索、差异路线
精确重组、全局最好精英重启、共享设施流交叉交换以及低频5路线邻域。第三轮进一步采用
停滞计数而非固定墙钟作为停止准则，并在检查点上重启。全局谱系401将17595 min降至
17409 min，谱系402继续降至17231 min；其中第124次迭代的普通低利用率4路线精确修复
完成$4\to3$吸收，一次节约100 min并把架次数由96降至95。随后定向5路线邻域贡献1 min，
另一条探索运行中的有效改进均来自普通incumbent序列，最终纯几何续算保持17218 min一级
目标，并把人员在途时间降至254656 min。故本文不把最终收益归因于机器学习。

强制source absorption的22个诊断窗口及前五个4/5/6路线窗口均未产生更优incumbent；
一般6路线局部主问题出现明显计算膨胀，予以弃用。路径重连V2四次尝试均未接受，交叉交换
在第三轮11次调用中没有一级收益；二者保留为实验对照，不进入最终精简主线。相反，
LAND-heavy算子在两条主谱系中累计贡献60 min，因此保留。

问题二完整流程为：生成燃油可行候选并得到联合RMP基准；以精确局部MILP依次执行固定
4路线搜索、长列增广、精英重组、全局重启、停滞检查点续算和低频5路线强化；始终保留词典序历史最优；最后按
区间着色回填人员编号，导出两个CSV并由独立Validator逐航段验收。

\subsection{算法实现与输出文件}

问题二复用\path{src/solver/evaluator.py}、
\path{technical\_stops.py}、\path{exporter.py}和独立Validator。
新增的\path{src/solver/q2.py}实现候选序列、取送兼容判定和全局RMP；
\path{src/solver/q2\_lns.py}实现破坏算子、有界长列生成、精确局部MILP、模拟退火和精英重组；
\path{src/solver/q2\_flow.py}构造有向OD流特征，\path{src/solver/q2\_round2.py}实现精英池、
差异路径重组与局部分支可行性门，\path{src/solver/q2\_round3.py}实现吸收潜力排序、
盆地指纹、局部主问题审计和有界深挖队列，\path{src/solver/q2\_artifacts.py}保证正式结果原子固化。
基础RMP入口为\path{scripts/05\_solve\_q2.py}，强化入口为
\path{scripts/06\_optimize\_q2\_lns.py}，精英重组入口为
\path{scripts/07\_recombine\_q2\_elites.py}；第三轮审计、吸收搜索和汇总入口分别为
\path{scripts/11\_audit\_q2\_round3.py}、\path{scripts/12\_search\_q2\_absorption.py}和
\path{scripts/13\_summarize\_q2\_round3.py}。图表由
\path{scripts/plot\_q2\_results.py}与\path{scripts/plot\_integrated\_improvements.py}
从最终JSON、比较表和CSV生成。正式目录仅保留同一复现运行的CSV、指标、配置和Validator，
大规模候选事件日志不进入支撑材料。

\subsection{结果展示与算法分析}

表~\ref{tab:q2-results}比较三类需求分开运输、候选主问题和精确局部搜索的主要阶段。
分开运输基线的三个子问题内部仍允许同类OD合并，仅禁止三类需求共享架次，因此不是
刻意构造的弱基线。三种方案均覆盖4000人并通过独立校验。

\begin{table}[htbp]
  \centering
  \caption{问题二不同算法的结果比较}
  \label{tab:q2-results}
  \resizebox{\linewidth}{!}{%
  \begin{tabular}{lrrrrrr}
    \toprule
    方法 & 飞机时间/min & 人员在途/min & 架次数 & 燃油/kg & 利用率/\% & 运行时间/s \\
    \midrule
    三类需求分开运输（B0） & 44184 & 260521 & 258 & 353993.8 & 34.661 & 310.857 \\
    最小联合候选池（B1）   & 20759 & 262954 & 114 & 162325.5 & 75.156 & 222.221 \\
    候选路线整数主问题（RMP） & 19736 & 270734 & 107 & 152910.4 & 81.830 & 301.609 \\
    几何长列精确LNS          & 19143 & 269218 & 103 & 148161.8 & 84.499 & -- \\
    第一轮强化解             & 17958 & 263588 & 97 & 138075.2 & 89.746 & -- \\
    第二轮多重启与路线消除   & 17595 & 259487 & 96 & 135954.1 & 91.021 & -- \\
    第三轮吸收式强化最终解   & \textbf{17218} & \textbf{254656} & \textbf{95}
                                  & \textbf{132473.9} & \textbf{92.703} & 2451.172$^{*}$ \\
    \bottomrule
  \end{tabular}}
\end{table}

联合运输相对分开运输基线的总飞机使用时间降低率为
\begin{equation}
  I_{\mathrm{joint}}
  =\frac{44184-17218}{44184}\times100\%=61.031\%.
  \label{eq:q2-joint-improvement}
\end{equation}
联合优化同时减少163个架次和221519.9 kg燃油，座位利用率提高58.042个百分点；人员
总在途时间减少5865 min，即2.251\%，因此五项指标均优于分开运输基线。相对旧RMP，
最终解再减少2518 min、12个架次和20436.5 kg燃油；相对第二轮再减少377 min并少1架次。
最终T1、T2、T3分别执行1、35、59架次，平均每架累计运输42.105人。表中带$^{*}$的时间为
获胜谱系各段累计到最终词典序最好值的时间，不等同于一次连续运行的墙钟时间。

\subsubsection{候选主问题求解界限}

设求解器在一级主问题上返回的当前最好整数解为$Z^{\mathrm{I}}$，对偶下界为
$Z^{\mathrm{LB}}$。本文按求解器定义计算候选集内相对间隙
\begin{equation}
  \operatorname{Gap}_{\mathrm{MIP}}
  =\frac{Z^{\mathrm{I}}-Z^{\mathrm{LB}}}{|Z^{\mathrm{I}}|}\times100\%.
  \label{eq:q2-mip-gap}
\end{equation}
求解界限见表~\ref{tab:q2-mip-bound}。旧双设施RMP的19736 min整数解与16296 min对偶界
对应17.430\%间隙；该证书只针对3116条固定候选路线。第二轮最终4路线修复的5.8651\%
间隙也只覆盖当次24条候选序列和1221个兼容分配。第三轮实现96到95架次吸收的4路线局部
主问题含806个兼容分配，并在该有限邻域内达到0间隙；这只能证明“给定被移除路线和本轮
候选列时无更好局部组合”，不能证明17218 min为完整Q2原问题全局最优。

\begin{table}[htbp]
  \centering
  \caption{问题二候选路线整数主问题的求解界限}
  \label{tab:q2-mip-bound}
  \begin{tabularx}{\linewidth}{XrrrrX}
    \toprule
    界限范围 & 路线/序列数 & 兼容分配数 & 整数值/min
    & 相对间隙/\% & 证明范围 \\
    \midrule
    历史双设施RMP & 3116 & 23251 & 19736 & 17.430 & 固定候选池 \\
    第二轮4路线局部修复 & 24 & 1221 & 局部目标 & 5.865 & 单次有限邻域 \\
    第三轮$4\to3$吸收邻域 & 有限4路线池 & 806 & 局部目标 & 0.000 & 单次有限邻域 \\
    完整Q2原问题 & -- & -- & 17218 & -- & 尚无全局证书 \\
    \bottomrule
  \end{tabularx}
\end{table}

需要强调的是，局部间隙小并不意味着全局间隙小；本文只把17218 min称为限定预算内的
当前最佳可行值，不声称Q2已达到全局最优。

图~\ref{fig:q2-convergence}展示从第一轮精确LNS到最终路线消除的改进链，飞机时间持续
下降，并在第三轮将架次数由96降至95。

\begin{figure}[htbp]
  \centering
  \IfFileExists{figures/q2_round3_improvement.pdf}
    {\includegraphics[width=0.94\linewidth]{figures/q2_round3_improvement.pdf}}
    {\figplaceholder{Q2三种求解阶段的飞机时间、架次数与利用率对比}}
  \caption{问题二精确局部重构与多重启的阶段改进}
  \label{fig:q2-convergence}
\end{figure}

选择一条同时运输出海、海返和穿梭人员的T3代表性架次，如图~\ref{fig:q2-route-state}。
该架次依次执行A03--F039--F048--F038--F045--A03，累计运输74人，其中出海、海返和
穿梭人员分别为19、18和37人；各航段离站载荷为19、18、0、19和18人，均不超过19座
容量，并在F038按精确燃油状态决定加油。
累计人数远大于座位数而逐航段合法，直观说明了先下后上和座位连续复用机制。

\begin{figure}[htbp]
  \centering
  \IfFileExists{figures/q2_representative_route.pdf}
    {\includegraphics[width=0.94\linewidth]{figures/q2_representative_route.pdf}}
    {\figplaceholder{Q2代表性混合架次：航段载荷、到达余油与加油位置}}
  \caption{问题二代表性混合架次的状态演化}
  \label{fig:q2-route-state}
\end{figure}

表~\ref{tab:q2-ablation}给出主要组件的逐层结果。扩展搜索与精英重组均被保留；专门的
局部分支和无界多样性伙伴因未带来可靠一级目标收益而不进入最终主线。

\begin{table}[htbp]
  \centering
  \caption{问题二精确局部搜索的逐层消融结果}
  \label{tab:q2-ablation}
  \begin{tabular}{lrrrrr}
    \toprule
    阶段 & 飞机时间/min & 人员在途/min & 架次数 & 燃油/kg & 是否采用 \\
    \midrule
    历史候选主问题 & 19736 & 270734 & 107 & 152910.4 & 基准 \\
    几何长列LNS & 19143 & 269218 & 103 & 148161.8 & 是 \\
    质量约束精英重组 & 17837 & 261114 & 97 & 138053.1 & 是 \\
    交叉交换 & 17693 & 258831 & 97 & 136752.4 & 是 \\
    第二轮路线消除 & 17595 & 259487 & 96 & 135954.1 & 是 \\
    第三轮吸收式强化 & \textbf{17218} & \textbf{254656} & \textbf{95} & \textbf{132473.9} & 是 \\
    \bottomrule
  \end{tabular}
\end{table}

最终\texttt{q2-validator.json}给出\texttt{valid=true}、
\texttt{issue\_count=0}；程序内存指标与从最终CSV复算的五项指标完全一致。全部4000人均
恰好分配一次、同架次完成取送且在首次到达终点时下机。完整95条路线、4000条人员记录、
最终复现配置和校验报告放入支撑材料；体量巨大的候选事件日志不随论文提交。

\section{问题三：带时间要求的多日排班}
\label{sec:q3}

\subsection{问题结构与总体思路}

问题三在问题二的4000条取送需求基础上增加了最早离开起点时刻、最晚送达时刻和
业务类型，并将飞机资源限制为24架具有固定所属机场和机型的具体直升机。规划期为
2026年8月3日00:00至8月10日00:00，共7个自然日。任一架次只能在06:00至18:00之间
起飞，必须在20:00前返回飞机所属机场；同一架飞机相邻两架次之间至少保留30 min
机场周转时间，且不得在海上设施过夜。

四类业务中，应急处置、增储上产和常规倒班人员共3840人，必须全部完成运输；160名
临时任务人员允许取消。题目要求先排除临时任务，求得完成全部非临时任务的基准总飞机
使用时间 $T_0$；随后在 $T_{\mathrm{air}}\le T_0$ 的条件下尽可能多地服务临时人员。
因此，问题三不是问题二路线的简单加时，而是带时间窗、固定属地异构机队和严格二阶段
目标的多日取送排班问题。

若直接联合优化4000名人员、路线、时刻、等待、加油和24架具体飞机，将产生大量组合
变量和对称解。本文采用“燃油可行路线复用--时间窗拼载--具体飞机区间排班--整数重分配
与多航班重构”的两层分解框架。路线层复用问题二的3116条燃油可行路线变体，并建立
23251个OD--路线兼容关系；排班层在人员可行起飞区间内逐架安排航班，用24个互不重叠
的飞机日历直接落实属地、机型和30 min周转限制。第一阶段完成全部非临时人员，第二
阶段先在固定架次上以稀疏0--1整数规划全局重分配人员，再允许同日两个相关架次重构为
一个新架次。每轮第二阶段解均删除临时人员并反馈为第一阶段候选，从而严格保证
$T_{\mathrm{air}}\le T_0$，且不会为保留临时任务而人为放宽第一阶段上界。

\subsection{问题三集合、变量与时间口径}

令规划日集合为 $\mathcal{D}=\{0,1,\ldots,6\}$，其中第 $q$ 日相对规划起点的分钟偏移
为 $\Omega_q=1440q$。非临时人员集合记为
$P_{\mathrm{req}}=P\setminus P_{\mathrm{temp}}$。路线层生成的候选架次集合记为
$\mathcal{C}$；对任意 $r\in\mathcal{C}$，已知其所属日期 $q(r)$、起降机场 $a(r)$、
机型 $k(r)$、停靠序列、加油选择以及候选人员集合 $P_r$。能够执行候选架次 $r$ 的
具体飞机集合为
\begin{equation}
  H_r=\{u\in H:a(u)=a(r),\ k(u)=k(r)\}.
  \label{eq:q3-compatible-aircraft}
\end{equation}
其中，$u$ 仅表示具体飞机索引，$h$ 仅表示路线停靠位置序号，与全文统一符号一致。

候选架次选择变量和飞机分配变量分别为
\begin{equation}
  x_r\in\{0,1\},\qquad y_{ru}\in\{0,1\},
  \quad r\in\mathcal{C},\ u\in H_r.
  \label{eq:q3-selection-variables}
\end{equation}
若选择候选架次 $r$，则 $x_r=1$；若该架次由飞机 $u$ 执行，则 $y_{ru}=1$。
记 $T_{rh}^{\mathrm{arr}}$ 和 $T_{rh}^{\mathrm{dep}}$ 分别为架次 $r$ 到达和离开第
$h$ 个停靠点的绝对分钟时刻。对于候选架次中的人员 $p\in P_r$，仍以
$\pi_{pr}$ 和 $\delta_{pr}$ 表示其上机和下机位置。临时人员服务变量沿用符号表中的
$z_p\in\{0,1\}$。

\subsection{时间窗感知的候选架次生成}

\subsubsection{紧任务排序与需求聚合}

问题三沿用前文任务排序指标
\begin{equation}
  \Pi_p=\bigl(\operatorname{prio}(c_p),l_p,S_p\bigr),
  \label{eq:q3-route-priority}
\end{equation}
并按照词典序升序选择待安排人员。程序分别试验“业务优先级--最晚到达--松弛量”、
“最晚到达--业务优先级--松弛量”和“松弛量--业务优先级--最晚到达”三种排序。
同一架次允许混合多个OD；装载时逐航段累计人数，因此人员下机后释放的座位可被后续
上机人员复用。业务优先级只决定构造顺序，不能替代硬约束，全部3840名非临时人员仍须
逐人覆盖。

\subsubsection{由问题二方案扩充候选池}

首先读取通过独立校验的问题二候选缓存，将每条路线的相对到离站时刻投影到7个规划日。
对人员$p$和候选路线$r$，由上、下机位置计算允许的共同起飞区间
\begin{equation}
  B_{prq}=\left[\max\{e_p-d_{r,\pi_{pr}},\Omega_q+360\},
  \min\{l_p-a_{r,\delta_{pr}},\Omega_q+1080,
  \Omega_q+1200-\tau_r\}\right].
  \label{eq:q3-start-interval}
\end{equation}
仅当该区间非空时保留兼容关系。为某个紧任务选择路线后，程序扫描所有尚未安排人员，
在共同起飞时刻上执行逐航段区间装箱；评分依次最小化单位人员飞机时间、架次人数的
相反数、架次时长和最早可行起飞时刻。这样既继承问题二路线的燃油安全性，又显式处理
问题三的时间窗和有限机队。

\subsection{候选架次的时刻递推与硬约束}

\subsubsection{逐站时刻与人员时间窗}

所有时刻均转化为自2026年8月3日00:00起的整数分钟。对于已选择的候选架次
$r$，相邻停靠点满足
\begin{equation}
  T_{r,h+1}^{\mathrm{arr}}
  =T_{rh}^{\mathrm{dep}}
  +t_{v_{rh},v_{r,h+1}}^{k(r)},
  \qquad h=0,1,\ldots,m_r.
  \label{eq:q3-time-recursion}
\end{equation}
中间海上停靠点的离开时刻满足
\begin{equation}
  T_{rh}^{\mathrm{dep}}
  \ge T_{rh}^{\mathrm{arr}}+10+10\xi_{rh},
  \qquad h=1,\ldots,m_r.
  \label{eq:q3-dwell-wait}
\end{equation}

程序为每条重构候选建立绝对整数分钟时刻表：对给定路线、人员和日期，从最晚可行起飞
时刻向前搜索，并逐站前推到达时刻、最短停留和人员最早上机约束；若必须等待，则把等待
量加入离站时刻。这样可识别“整体平移不可行、在中间设施等待后可行”的组合。最终被接受
的两个2--1合并无需额外等待，但时间感知模块扩大了邻域可行域；所有等待若出现均同时计入
飞机时间和相应人员在途时间，导出CSV时使用架次自身的实际到离站时刻。

若 $p\in P_r$ 且 $x_r=1$，则人员时间窗约束为
\begin{equation}
  T_{r,\pi_{pr}}^{\mathrm{dep}}\ge e_p,
  \qquad
  T_{r,\delta_{pr}}^{\mathrm{arr}}\le l_p.
  \label{eq:q3-passenger-window}
\end{equation}
人员 $p$ 的实际在途时间为
\begin{equation}
  T_{pr}^{\mathrm{pass}}
  =T_{r,\delta_{pr}}^{\mathrm{arr}}
  -T_{r,\pi_{pr}}^{\mathrm{dep}}.
  \label{eq:q3-passenger-time}
\end{equation}

\subsubsection{机场运营时段与同日闭合}

对于 $q=q(r)$ 且 $x_r=1$，候选架次必须满足
\begin{equation}
  \Omega_q+360
  \le T_{r0}^{\mathrm{dep}}
  \le \Omega_q+1080,
  \qquad
  T_{r,m_r+1}^{\mathrm{arr}}\le\Omega_q+1200.
  \label{eq:q3-airport-hours}
\end{equation}
其中360、1080和1200分别对应06:00、18:00和20:00。候选路线首末节点均为飞机
所属机场，且式~\eqref{eq:q3-airport-hours}保证架次在同一运营日内完成，因此不会在海上
设施过夜。

\subsubsection{具体飞机分配与周转约束}

每个被选候选架次恰好分配给一架兼容飞机：
\begin{equation}
  \sum_{u\in H_r}y_{ru}=x_r,
  \qquad r\in\mathcal{C}.
  \label{eq:q3-aircraft-assignment}
\end{equation}
对每个已分配的 $(r,u)$ 建立占用区间
\begin{equation}
  I_{ru}=\left[
    T_{r0}^{\mathrm{dep}},
    T_{r,m_r+1}^{\mathrm{arr}}+30
  \right),
  \quad y_{ru}=1.
  \label{eq:q3-optional-interval}
\end{equation}
并对每架具体飞机施加区间互斥条件
\begin{equation}
  \operatorname{NoOverlap}\bigl(\{I_{ru}:r\in\mathcal{C},\ u\in H_r\}\bigr),
  \qquad u\in H.
  \label{eq:q3-no-overlap}
\end{equation}
程序使用按起飞时刻排序的日历插入实现式~\eqref{eq:q3-no-overlap}：区间末端附加30 min，
使同一飞机下一架次的起飞时刻不早于上一架次返场后30 min。
最终的 \texttt{flight\_no} 不作为优化变量，而是在求解结束后按每架飞机的实际起飞时刻
排序并从1开始连续生成。

路线层已经验证同机场闭合、最多5次海上着陆、首次到达下机、逐航段容量和燃油安全；
排班层改变等待或起飞时刻后，仍须由独立仿真器使用
$\nu_k,\beta_k,G_k^{\max},G_k^{\min}$ 重新逐航段验证，避免候选属性与最终CSV不一致。

\subsection{第一阶段：非临时任务基准优化}

对每名必运人员定义其可服务候选集合
$\mathcal{C}_p=\{r\in\mathcal{C}:p\in P_r\}$。第一阶段人员覆盖约束为
\begin{equation}
  \sum_{r\in\mathcal{C}_p}x_r=1,
  \qquad p\in P_{\mathrm{req}}.
  \label{eq:q3-mandatory-cover}
\end{equation}
若 $x_r=1$，定义候选架次的实际飞机使用时间为
\begin{equation}
  T_r^{\mathrm{air}}
  =T_{r,m_r+1}^{\mathrm{arr}}-T_{r0}^{\mathrm{dep}};
  \label{eq:q3-flight-duration}
\end{equation}
若 $x_r=0$，则令 $T_r^{\mathrm{air}}=0$。第一阶段按照
\begin{equation}
  \min_{\mathrm{lex}}
  \left(
    T_{\mathrm{air}},
    T_{\mathrm{passenger}},
    N_{\mathrm{flight}},
    F_{\mathrm{fuel}}
  \right),
  \qquad
  T_{\mathrm{air}}=\sum_{r\in\mathcal{C}}T_r^{\mathrm{air}},
  \label{eq:q3-stage1-objective}
\end{equation}
进行多阶段词典序求解。实际实现时先最小化 $T_{\mathrm{air}}$；得到限定时间内的最佳
可行值后固定该值，再依次优化其余指标，避免任意加权导致后级指标抵消一级飞机时间。
将通过独立校验器的第一阶段方案记为 $S_{\mathrm{base}}$，并定义
\begin{equation}
  T_0=T_{\mathrm{air}}(S_{\mathrm{base}}).
  \label{eq:q3-baseline-time}
\end{equation}
若求解器未证明全局最优，论文应称 $T_0$ 为“限定时间内最佳可行基准值”，不能直接称为
理论最优值。

实际求解采用多启发式启动与Top-$K$筛选。除业务优先、最晚到达和最小松弛量外，程序预先
计算每名人员的可行日期数、兼容路线数、最短与次短路线差、OD需求密度和时间窗宽度，形成
日期稀缺、路线稀缺、关键度、OD密度、Regret代理及Hard-Skeleton/Flexible-Regret等11种
确定性构造模式，并增加1个固定种子的随机化启动，共形成12个启动。先对全部模式执行低成本完整构造，仅对词典序最优且结构有差异的Top-$K$候选
进行整数重分配和深度邻域。若某任务找不到日历空隙，程序退化为单人架次并扫描全部兼容
路线、日期和飞机；仍失败时记录冲突人员。固定随机种子的模式可完全复现。
\subsection{第二阶段：临时任务最大满足优化}

第二阶段保留全部非临时人员覆盖约束，并对临时人员设置
\begin{equation}
  \sum_{r\in\mathcal{C}_p}x_r=z_p,
  \qquad p\in P_{\mathrm{temp}}.
  \label{eq:q3-temporary-cover}
\end{equation}
首先尝试把临时人员插入起终点已经按正确顺序出现、相应航段仍有空座且时间窗可行的
已选路线；此类插入不改变停靠序列和时刻，因而不会增加总飞机使用时间。随后固定当前
架次的路线、时刻和具体飞机，对全部4000人重新建立人员--架次0--1分配变量$w_{pr}$。
每名必运人员满足$\sum_r w_{pr}=1$，临时人员满足$\sum_r w_{pr}\le1$，并对每个架次的
每一航段施加$\sum_{p:\pi_{pr}\le h<\delta_{pr}}w_{pr}\le s_{k(r)}$。最终固定架次模型含
8630个0--1变量和4575条约束，始终施加
\begin{equation}
  T_{\mathrm{air}}\le T_0.
  \label{eq:q3-time-upper-bound}
\end{equation}
第二阶段的首要目标为
\begin{equation}
  \max \sum_{p\in P_{\mathrm{temp}}}z_p.
  \label{eq:q3-stage2-objective}
\end{equation}
先最大化临时人员服务数；固定该整数最优值后，再最小化人员总在途时间。因此在固定架次
集合内执行严格词典序，而不是使用可能发生量纲抵消的加权和。若进一步允许改变路线，
可再按照
\begin{equation}
  \min_{\mathrm{lex}}
  \left(
    T_{\mathrm{air}},
    T_{\mathrm{passenger}},
    N_{\mathrm{flight}},
    F_{\mathrm{fuel}},
    -U
  \right)
  \label{eq:q3-stage2-secondary}
\end{equation}
优化其他指标。临时任务满足率为
\begin{equation}
  R_{\mathrm{temp}}
  =\frac{\sum_{p\in P_{\mathrm{temp}}}z_p}{160}\times100\%.
  \label{eq:q3-temp-rate}
\end{equation}
未安排的临时人员仍须保留在 \texttt{q3-assignments.csv} 中，其架次、飞机和时间字段
按题目要求留空；3840名非临时人员的安排字段不得为空。

\subsubsection{两阶段投影反馈与多航班结构邻域}

设第二阶段可行方案为$S_2$，定义投影算子$\mathcal{P}_{\mathrm{req}}(S_2)$：保留全部
航班、路线、具体飞机和时刻，仅删除临时人员的分配。删除乘客不会增加航段载荷、油耗、
飞行时间或时间窗约束，因此
\begin{equation}
  T_{\mathrm{air}}\!\left(\mathcal{P}_{\mathrm{req}}(S_2)\right)
  =T_{\mathrm{air}}(S_2),
  \qquad
  \mathcal{P}_{\mathrm{req}}(S_2)\ \text{为第一阶段可行解}.
  \label{eq:q3-stage2-projection}
\end{equation}
程序据此执行有限反馈闭环：第一阶段解确定上界后求第二阶段，再把第二阶段投影回第一
阶段并执行仅含非临时人员的深度优化；若飞机时间严格下降，则用新上界重求第二阶段，
否则终止。该逻辑纠正了旧方案中“第二阶段30510 min反而小于第一阶段30546 min、却未
反馈更新”的不一致。

为突破逐架替换的局部最优，对同日航班按设施Jaccard相似度、所属机场和起飞时间构造
近邻图，从中选取2--4架相关航班进行Ruin。程序在3116条缓存路线中构造可变数量的替代架次，联合选择路线、机型、具体飞机和时间感知时刻表，再用全局0--1分配模型修复全部人员。仅接受必运人员完整、临时服务数不下降、飞机日历合法且符合对应阶段词典序的方案。Closure阶段找到一次真实$3\to2$重构，节约211 min并将架次由171降至170。

通过Closure Gate后，P2首先对未服务临时人员自动生成“路线--时间--容量--日历”可行性档案，再按固定架次、2架、3架和4--5架相关航班逐级恢复；随后使用相关航班、高成本、低利用率、飞机链和随机相关五类破坏算子执行自适应大邻域。对宽时间窗换班人员，跨日邻域允许删除低装载架次并在其他日已有架次中重分配。本次P2删除两个冗余架次，经最终第二阶段投影反馈后得到168架次、29659 min的统一结构。

\subsubsection{PRO V1精英池邻域与PRO V2深度强化}

Closure--P2收敛后，程序把历代通过校验的高质量解（含29659 min、30180 min等12个精英结构）固化为精英池，PRO V1在其上执行路线库定价、精英间结构距离度量和受控破坏--修复，把第一阶段进一步降至29155 min（165架次）。PRO V2随后从三个方向继续深挖。

其一是异构参数筛选：以20组差异化的算子概率、重启阈值、分组规模和路线限额配置各配2个固定种子，共40次受预算保护的短程运行，分别收敛到29155、29006和28868 min三档（22、13、5次），说明单一参数盆地之外仍存在更优结构。其二是多岛深度ALNS：将筛选得到的代表性配置分为抛光、跨日、多样性和飞机链4个岛，每岛最多500次迭代、900 s墙钟保护，四岛均收敛到28868 min。表~\ref{tab:q3-operator-v2}汇总筛选与深度两阶段合计的算子统计：766次算子尝试中接受29次、产生23次全局最好，其中跨日算子占28次接受和全部23次全局最好，是29155$\to$28868这287 min改进的主导结构来源；这也解释了为何同日局部修复在该盆地收益枯竭。

\begin{table}[htbp]
  \centering
  \caption{问题三PRO V2筛选与深度阶段的算子贡献统计}
  \label{tab:q3-operator-v2}
  \begin{tabular}{lrrr}
    \toprule
    算子类别 & 尝试次数 & 接受次数 & 全局最好次数 \\
    \midrule
    跨日结构搜索 & 128 & 28 & 23 \\
    路线抛光与其他同日算子 & 638 & 1 & 0 \\
    \midrule
    合计 & 766 & 29 & 23 \\
    \bottomrule
  \end{tabular}
\end{table}

其三是Optional Rescue与强制P0反馈：对28868 min结构中仍未服务的临时人员构造深度2--4、组规模至多12的定向救援邻域，接受1次结构移动；随后把第二阶段删除临时人员的投影作为强制P0反馈回第一阶段再优化，最终把飞机时间降至28728 min（162架次），第二阶段临时服务数保持157人不变，即改进完全来自必运结构收紧而非放宽临时目标。为检验局部最优性，程序还在28868 min解上执行了三类精确邻域：50个引导式精确LNS窗口（合计1876.9 s）、半径5/10/20/40/80的局部分支和10个飞机--日链窗口，均未接受任何改进；这些是有限邻域内的负证据，不构成全局最优性证明。定价阶段另把363条新选路线导入主池，最终持久化资产为12个精英（平均成对结构距离0.0519）、3116条路线变体和来自12个排班来源的179条去重航班列。

\subsection{问题三求解界限}

为避免把启发式可行值误称为全局最优，本文把第一阶段界限分为“对原问题有效的全局
下界”和“仅对有限候选路线池有效的参考界”两类，并另行给出第二阶段人数目标的上界。

\subsubsection{客公里工作量下界}

令$d_p^{\mathrm{sp}}$为人员$p$起终点之间的最短路距离；当端点为LAND时，对三个机场取
最小值。任一时刻全机队把客公里转化为飞机时间的效率不可能超过单架飞机的最大
座位公里率
\begin{equation}
  \rho_{\max}=\max_{k\in K}\frac{s_k\nu_k}{60}.
  \label{eq:q3-seat-km-rate}
\end{equation}
故第一阶段存在对原问题有效的运输工作量下界
\begin{equation}
  LB_3^{\mathrm{air}}
  =\left\lceil
  \frac{\sum_{p\in P_{\mathrm{req}}}d_p^{\mathrm{sp}}}{\rho_{\max}}
  \right\rceil=12389\ \mathrm{min}.
  \label{eq:q3-aircraft-time-lb}
\end{equation}
该下界乐观地忽略返场空驶、停靠、加油绕行、时间窗和具体飞机冲突，所以数值偏松，
但不会高估原问题最优值。

\subsubsection{机场--机型分层连续多商品流下界}

为把返场结构、真实航段时间、海上停靠时间、机型座位数和属地机队总工时纳入下界，
构造机场--机型层集合$\mathcal{L}=A\times K$。对层$\ell=(a,k)$，令$\mathcal{E}_\ell$为以机场$a$为基地产生
的机型$k$可飞有向弧集，连续变量$X_{ij}^{\ell}$表示该层在弧$(i,j)$上的飞机流量。弧成本取
\begin{equation}
  c_{ij}^{\ell}=\left\lceil\frac{60d_{ij}}{\nu_k}\right\rceil
  +10\,\mathbf{1}_{\{j\in F\}},
  \qquad \ell=(a,k),
  \label{eq:q3-layered-arc-cost}
\end{equation}
其中第二项计入到达海上设施后的最小停靠时间。每个层满足飞机流守恒
\begin{equation}
  \sum_{j:(i,j)\in \mathcal{E}_\ell}X_{ij}^{\ell}
  -\sum_{j:(j,i)\in \mathcal{E}_\ell}X_{ji}^{\ell}=0,
  \qquad i\in V,\ \ell\in \mathcal{L}.
  \label{eq:q3-layered-aircraft-balance}
\end{equation}

为控制模型规模，将旅客按目的地$d$聚合成连续商品。令$F_{ij}^{d}$为目的地为$d$的
旅客在弧$(i,j)$上的流量，$b_i^d$为节点$i$对该商品的净供给；LAND通过虚拟汇点与三个
机场相连，从而保留“可在任一机场完成陆地端运输”的题意。旅客流守恒和航段容量分别为
\begin{align}
  \sum_jF_{ij}^{d}-\sum_jF_{ji}^{d}&=b_i^d,
  \label{eq:q3-passenger-flow-balance}\\
  \sum_dF_{ij}^{d}&\leq
  \sum_{\ell=(a,k)\in \mathcal{L}}s_kX_{ij}^{\ell}.
  \label{eq:q3-layered-seat-capacity}
\end{align}
同时用属地机队总可用时间限制各层飞机流：
\begin{equation}
  \sum_{(i,j)\in \mathcal{E}_\ell}c_{ij}^{\ell}X_{ij}^{\ell}
  \leq n_{ak}\times7\times840,
  \qquad \ell=(a,k),
  \label{eq:q3-layered-fleet-time}
\end{equation}
其中$n_{ak}$为机场$a$所属机型$k$的飞机数。连续松弛的目标为
\begin{equation}
  LB_3^{\mathrm{MCF}}
  =\min\sum_{\ell\in \mathcal{L}}\sum_{(i,j)\in \mathcal{E}_\ell}
  c_{ij}^{\ell}X_{ij}^{\ell}.
  \label{eq:q3-mcf-lower-bound}
\end{equation}

任一原问题可行排班都可按具体飞机的属地和机型映射为$X$，并按每名旅客的实际航段
叠加为$F$；反之，本模型允许连续架次、旅客拆分和换乘，并放松日期分配、时间窗、逐机
不重叠、30 min周转、加油、整数座位和最多经停数等约束。因此它的可行域包含原问题
所有可行解，目标值必不超过原问题最优值，是对原问题有效的全局下界。模型含
191124个连续变量、3613个等式和2973个不等式，HiGHS求得连续最优值
$14124.895$ min，向上取整后为14125 min。于是
\begin{equation}
  LB_3^{\mathrm{global}}
  =\max\left\{LB_3^{\mathrm{air}},\left\lceil LB_3^{\mathrm{MCF}}\right\rceil\right\}
  =14125\ \mathrm{min}.
  \label{eq:q3-enhanced-global-lb}
\end{equation}

\subsubsection{有限候选路线池LP参考界}

进一步把Q2生成的3116条燃油可行路线变体放入连续主问题：令$n_r^{\mathrm{LP}}\geq0$为路线$r$的
连续执行次数，$w_{gr}^{\mathrm{LP}}\geq0$为OD组$g$分配给路线$r$的人数，约束每组需求全部满足且
每条路线的每一航段均不超载。该模型的连续最优值为15197.678 min，向上取整为15198 min。
PRO V2定价阶段又在缩减后的229条候选变体、107条入选列上求解受限主问题，连续最优值
17217.417 min（上取整17218 min）。两个模型的共同点是都排除了池外路线，会抬高目标；
连续化和放松逐机排班又会压低目标，二者方向相反，故15198 min与17218 min都只能作为
“候选池内参考界”，不能作为原问题的全局最优性证书。

\begin{table}[htbp]
  \centering
  \caption{问题三第一阶段求解界限对照（相对可行上界28728 min）}
  \label{tab:q3-bound-comparison}
  \begin{tabularx}{\linewidth}{XrrX}
    \toprule
    界限模型 & 界限/min & 相对区间宽度 & 适用范围 \\
    \midrule
    客公里工作量下界 & 12389 & 56.875\% & 原问题全局有效 \\
    机场--机型分层连续多商品流 & 14125 & 50.832\% & 原问题全局有效 \\
    3116条路线有限池连续主问题 & 15198 & 47.098\% & 仅作候选池参考 \\
    定价缩减池受限主问题 & 17218 & 40.067\% & 仅作候选池参考 \\
    \bottomrule
  \end{tabularx}
\end{table}

本文可行上界为$UB_3^{\mathrm{air}}=T_0=28728$ min，采用增强全局下界后的认证区间
宽度为
\begin{equation}
  g_3^{\mathrm{air}}
  =\frac{UB_3^{\mathrm{air}}-LB_3^{\mathrm{global}}}{UB_3^{\mathrm{air}}}
  \times100\%=50.832\%.
  \label{eq:q3-aircraft-bound-gap}
\end{equation}
与客公里下界相比，分层流下界提高1736 min；相较Closure--P2的29659 min可行值，PRO V1/V2
把上界再降低931 min，使增强界的认证区间由52.375\%缩小到50.832\%。由于区间仍然
较宽，28728 min应表述为“当前最佳可行基准”，不能声称已证明第一阶段全局最优。

第二阶段的原问题上界直接来自临时人员总数：
\begin{equation}
  Z_{\mathrm{temp}}^*\le |P_{\mathrm{temp}}|=160.
  \label{eq:q3-temporary-upper-bound}
\end{equation}
在新上界$T_0=28728$ min下，最终方案服务157人，尚有P1102、P2239、P3290三人未服务，
相对理论上界的绝对差为3人、上界口径相对差为$3/160=1.875\%$。固定最终162个架次时，
人员分配MILP返回最优状态并证明157人为该固定结构的最优服务数；但允许改变路线和航班
结构的原问题仍可能达到158、159或160，三个取值均未关闭，因此本文不声称第二阶段全局
最优。旧方案在较宽上界30510 min下达到160人、29659 min上界下达到158人，均不能与
28728 min上界下的157人直接比较：由于第一阶段首要目标收紧，新157人结果在词典序意义
上严格优于旧158人结果，不构成退步，也不能据此放宽第一阶段上界。

\subsection{冲突反馈的两层迭代算法}

问题三完整求解流程如下：
\begin{enumerate}[label=\textbf{步骤\arabic*.},leftmargin=4.5em]
  \item 读取Q3人员、24架飞机和统一参数，转化7天内的绝对整数分钟；
  \item 计算每名人员的乐观最短运输时间、可行日期数、兼容路线数、Regret代理、OD密度
  与时间窗宽度，形成11种可复现构造键；
  \item 将3840名非临时人员设为必运集合，将160名临时人员设为第二阶段可选集合；
  \item 读取Q2的3116条燃油可行路线变体，建立OD--路线兼容索引；
  \item 求解机场--机型分层连续多商品流松弛和有限候选池连续主问题，分别记录原问题
  全局下界与候选池参考界；
  \item 由式~\eqref{eq:q3-start-interval}生成个人--路线--日期可行起飞区间；
  \item 选取最紧人员作为种子，在共同起飞时刻逐航段装填其他兼容人员；
  \item 在所属机场和机型对应的飞机日历中寻找最早空隙，并预留30 min周转；
  \item 若找不到群组方案，则扫描单人候选并记录冲突人员，不删除已有可行架次；
  \item 对12个启动分别构造方案，先低成本筛选，再对Top-$3$多样候选执行深度优化；
  \item 固定架次建立稀疏0--1人员重分配模型，使人员组合不受构造顺序锁定；
  \item 在保持人员集合可行的条件下，以更短的同机场、同机型候选替换冗余路线；
  \item 允许整架次更换机型、所属机场、具体飞机和起飞时刻，并重新检查时间窗与日历冲突；
  \item 执行删停靠--全局人员整数重分配，以及同日2--4架旧航班到可变替代数的
  $k\to m$ Ruin--Recreate；候选时刻表允许海上等待并使用实际到离站时刻；
  \item 将旧第二阶段解删除临时人员后正式投影为第一阶段候选，Closure中以$3\to2$重构得到29969 min；
  \item 通过Closure Gate后，依次执行未服务临时人员定向恢复、五类自适应大邻域、动态路线增广和跨日柔性需求搜索，经投影反馈得到29659 min；
  \item PRO V1固化12个历史精英，执行路线库定价与受控破坏--修复，降至29155 min；
  \item PRO V2执行20配置$\times$2种子异构参数筛选与4岛深度ALNS，跨日结构搜索主导降至28868 min；
  \item 以Optional Rescue与强制P0反馈投影再优化，最终得到$T_0=28728$ min；
  \item 在新上界下最大化临时人员数，再复核所得解投影回第一阶段后上界不变；
  \item 导出162条架次和4000条人员记录，其中3名未服务临时人员保留空分配字段，用独立
  校验器复算全部硬约束和目标指标。
\end{enumerate}

\subsection{结果展示与算法分析}

表~\ref{tab:q3-stage1-comparison}给出可验证的逐层改进。旧深度解为30546 min；将旧第二
阶段解严格投影后，第一阶段立即降至30510 min。P0--P1闭环将其降至30180 min；Closure中的$3\to2$重构进一步降至29969 min，P2跨日柔性需求搜索与投影反馈降至29659 min。PRO V1精英池强化降至29155 min；PRO V2经异构筛选与多岛深度ALNS降至28868 min，再以Optional Rescue与强制P0反馈降至28728 min。相较旧深度解累计减少1818 min、11个架次和14772.4 kg燃油。由于一级目标优先，人员在途时间的小幅增加不构成退化。

\begin{table}[htbp]
  \centering
  \caption{问题三第一阶段算法对比}
  \label{tab:q3-stage1-comparison}
  \begin{tabularx}{\linewidth}{Xrrrrr}
    \toprule
    方法 & 可行性 & 飞机时间/min & 人员在途/min & 架次数 & 燃油/kg \\
    \midrule
    旧版交替搜索基准 & 是 & 30546 & 241130 & 173 & 241607.7 \\
    第二阶段删除临时人员投影 & 是 & 30510 & 240522 & 173 & 241133.6 \\
    P0--P1反馈与结构调整 & 是 & 30180 & 240363 & 171 & 239702.3 \\
    Closure：同日$3\to2$重构 & 是 & 29969 & 240594 & 170 & 238247.2 \\
    Closure--P2与最终反馈 & 是 & 29659 & 241073 & 168 & 235329.4 \\
    PRO V1精英池强化 & 是 & 29155 & 241018 & 165 & 231384.5 \\
    PRO V2筛选、深度与最终反馈 & 是 & \textbf{28728} & \textbf{242455} & \textbf{162} & \textbf{226835.3} \\
    \bottomrule
  \end{tabularx}
\end{table}

表~\ref{tab:q3-evolution}单独汇总PRO阶段之后的算法演进链：29155$\to$28868的287 min来自异构筛选与多岛深度ALNS（跨日算子主导），28868$\to$28728的140 min来自Optional Rescue与强制P0反馈，且第二阶段临时服务数全程保持157人。

\begin{table}[htbp]
  \centering
  \caption{问题三PRO V1/V2算法演进链}
  \label{tab:q3-evolution}
  \begin{tabularx}{\linewidth}{XrrX}
    \toprule
    阶段 & 飞机时间/min & 架次 & 主要改进来源 \\
    \midrule
    Closure--P2反馈 & 29659 & 168 & 跨日柔性需求搜索与投影反馈 \\
    PRO V1 & 29155 & 165 & 精英池路线定价与受控破坏--修复 \\
    PRO V2深度搜索 & 28868 & 163 & 异构参数筛选与4岛深度ALNS（跨日主导） \\
    最终解 & \textbf{28728} & \textbf{162} & Optional Rescue与强制P0反馈 \\
    \bottomrule
  \end{tabularx}
\end{table}

两阶段最终结果按表~\ref{tab:q3-results}汇总。在严格新上界28728 min下，固定架次人员
分配MILP服务157名临时人员，尚有P1102、P2239、P3290三人未服务。最终反馈把第二阶段航班结构
严格投影为第一阶段候选，因而两阶段的飞机时间、架次数和燃油完全一致；新增157名乘客使
人员总在途时间增加9322 min。按已服务人数平均，在途时间由63.139 min/人降至62.992 min/人。

\begin{table}[htbp]
  \centering
  \caption{问题三两阶段优化结果}
  \label{tab:q3-results}
  \begin{tabularx}{\linewidth}{Xrrrrrr}
    \toprule
    阶段 & 临时人数 & 飞机时间/min & 人员在途/min & 架次 & 燃油/kg & 利用率/\% \\
    \midrule
    非临时基准 $S_{\mathrm{base}}$ & 0 & 28728 & 242455 & 162 & 226835.3 & 57.840 \\
    第二阶段最终方案 & \textbf{157} & \textbf{28728} & \textbf{251777} & \textbf{162} & \textbf{226835.3} & \textbf{60.098} \\
    \bottomrule
  \end{tabularx}
\end{table}

图~\ref{fig:q3-gantt}以横轴表示7天绝对时间、纵轴表示具体飞机，颜色区分任务类型或
架次所属机场；每个矩形表示一架次，矩形之间的空白应能反映30 min周转。

\begin{figure}[htbp]
  \centering
  \IfFileExists{figures/q3_gantt.pdf}
    {\includegraphics[width=0.98\linewidth]{figures/q3_gantt.pdf}}
    {\figplaceholder[0.94\linewidth]{7天具体飞机排班甘特图：飞机ID、架次区间、30 min周转和任务类型}}
  \caption{问题三具体飞机的7天排班甘特图}
  \label{fig:q3-gantt}
\end{figure}

为验证人员时间窗，定义已服务人员的送达裕度
$M_p=l_p-T_{r,\delta_{pr}}^{\mathrm{arr}}$。图~\ref{fig:q3-time-margin}按业务类型展示
送达裕度分布，并单独标出裕度最小的紧任务。全部3997名已服务人员均满足 $M_p\ge0$，正文
复算得到最小裕度为0 min、迟到人数为0、全体裕度中位数1437 min；应急处置与增储上产的
裕度中位数分别为19 min和54.5 min，说明紧任务的可行性不是由宽时间窗任务掩盖的。

\begin{figure}[htbp]
  \centering
  \IfFileExists{figures/q3_time_window_margin.pdf}
    {\includegraphics[width=0.84\linewidth]{figures/q3_time_window_margin.pdf}}
    {\figplaceholder{不同业务类型的实际送达裕度分布及最紧任务标注}}
  \caption{问题三已服务人员的送达时间裕度}
  \label{fig:q3-time-margin}
\end{figure}

图~\ref{fig:q3-aircraft-utilization}以日期和具体飞机为两个维度，颜色表示每日飞机使用
分钟数或占可运营时长的比例，用于识别机场--机型资源瓶颈。该图必须按飞机所属机场分组，
不能把不同机场的同型飞机视为可互换资源。最高单日占用为A03-T2-H01在8月8日的
575 min，占06:00--20:00可用区间的68.452\%，不存在超过运营时段的隐性排班。

\begin{figure}[htbp]
  \centering
  \IfFileExists{figures/q3_aircraft_utilization.pdf}
    {\includegraphics[width=0.90\linewidth]{figures/q3_aircraft_utilization.pdf}}
    {\figplaceholder{24架飞机按7天统计的使用时间或利用率热力图}}
  \caption{问题三具体飞机的分日利用率热力图}
  \label{fig:q3-aircraft-utilization}
\end{figure}

表~\ref{tab:q3-ablation}给出从P0--P1到PRO V2的关键检查点。Closure在真实飞机日历
上完成$3\to2$重构；P2A定向恢复、P2B自适应大邻域和P2C动态路线增广在本次受限预算内
保持当前解，P2D跨日柔性搜索分别删除第一、二阶段的一条低效架次。PRO V1精英池强化与
PRO V2筛选/深度搜索继续压低上界，最后必须将第二阶段
删除临时人员后投影回第一阶段，才形成两阶段共用的28728 min上界。旧160人结果对应
30510 min上界、旧158人结果对应29659 min上界，均不能脱离第一阶段上界与新157人结果
直接比较。

\begin{table}[htbp]
  \centering
  \caption{问题三反馈优化与PRO强化的逐层消融结果}
  \label{tab:q3-ablation}
  \begin{tabularx}{\linewidth}{Xrrrr}
    \toprule
    方案 & 第一阶段/min & 架次 & 第二阶段临时人数 & 第二阶段/min \\
    \midrule
    旧版深度方案 & 30546 & 173 & 160 & 30510 \\
    P0：第二阶段严格投影 & 30510 & 173 & 160 & 30510 \\
    P0--P1闭环基准 & 30180 & 171 & 158 & 30180 \\
    Closure：同日$3\to2$重构 & 29969 & 170 & 158 & 29969 \\
    P2A--P2C：定向恢复、ALNS与增广 & 29969 & 170 & 158 & 29969 \\
    P2D：跨日柔性需求搜索 & 29853 & 169 & 158 & 29659 \\
    Closure--P2最终投影反馈 & 29659 & 168 & 158 & 29659 \\
    PRO V1精英池强化 & 29155 & 165 & 157 & 29155 \\
    PRO V2筛选与4岛深度ALNS & 28868 & 163 & 157 & 28868 \\
    Rescue与强制P0反馈最终解 & \textbf{28728} & \textbf{162} & \textbf{157} & \textbf{28728} \\
    \bottomrule
  \end{tabularx}
\end{table}

最后，表~\ref{tab:q3-stability}汇总12个启动的筛选结果，其中9个得到完整可行解；路线
稀缺、Regret代理以及完整硬骨架--柔性Regret构造落入不可行盆地，仅作为诊断而不进入
Top-$K$。可行初解飞机时间标准差为940.216 min，说明初始盆地影响明显。深度搜索始终把
28728 min现有可行解作为保护性候选，因此扩展多启动不会导致结果回归。

\begin{table}[htbp]
  \centering
  \caption{问题三12个启动中9个可行初解的稳定性结果}
  \label{tab:q3-stability}
  \begin{tabular}{lrrr}
    \toprule
    指标 & 最佳值 & 均值 & 标准差 \\
    \midrule
    初始构造飞机时间/min & 31371 & 32262.778 & 940.216 \\
    非临时人员总在途/min & 265248 & 267333.778 & 1448.039 \\
    架次数 & 173 & 179.222 & 5.391 \\
    运行时间/s & 40.541 & 61.648 & 13.363 \\
    \bottomrule
  \end{tabular}
\end{table}

完整的24架飞机时刻表、4000条人员记录、投影审计、反馈轨迹、多启动与异构筛选结果、
基准方案、最终方案和两份独立校验报告均放入支撑材料；最终CSV中3997条记录具有有效
分配，3条未服务临时任务按题意留空，两阶段Validator报告违规数均为0。
图~\ref{fig:q3-search-evolution}汇总上述算法演进的飞机时间改进链。

\begin{figure}[htbp]
  \centering
  \IfFileExists{figures/q3_search_evolution.pdf}
    {\includegraphics[width=0.86\linewidth]{figures/q3_search_evolution.pdf}}
    {\figplaceholder{问题三算法演进的飞机时间改进链}}
  \caption{问题三算法演进的飞机时间改进链（30546$\to$29659$\to$29155$\to$28868$\to$28728 min）}
  \label{fig:q3-search-evolution}
\end{figure}

\section{模型评价与推广}
\subsection{模型优点}
\begin{enumerate}[label=\arabic*.]
  \item 以统一状态递推处理三问，指标和硬约束口径一致；
  \item OD聚合、Q1的ALNS与Q2的候选路线MILP兼顾4000人规模和复杂取送约束；
  \item 路线层与区间日历排班层分解，并以稀疏整数规划修复人员分配，兼顾可行性与可复现性；
  \item 独立校验器从最终 CSV 复算，使模型、实现、论文和提交文件形成闭环。
\end{enumerate}

\subsection{模型局限}
模型采用确定性距离、飞行时间和油耗，未考虑天气、临时禁飞、停机坪并发及故障等随机扰动；Q1虽以路线池主问题把可行值降至14730 min，但相对原问题全局下界的保守区间仍为10.445\%，14208.637 min只是不完备路线池LP参考界。Q2的16296 min对偶界属于旧有限候选RMP，第三轮吸收邻域的0间隙也只覆盖一次有限4路线局部主问题，故17218 min尚无原问题全局证书；机器学习数据虽已达到离线建模门槛，本轮没有训练模型，也不把探索标签当作收益原因。Q3以机场--机型分层多商品流把第一阶段全局下界提高到14125 min，但28728 min可行值与该下界之间仍有50.832\%的保守区间，15197.6776 min与17217.4167 min分别仅为3116候选池LP和定价缩减池受限主问题的池内参考界。第二阶段157人只在固定162架次的分配模型中得到最优证明，距离原问题人数上界160仍差3人，且不限定架次结构的158、159与160人可行性均未关闭。后续可为Q1路线池主问题增加对称性消除和MIP start，对Q2开展按谱系分组的LR排序A/B，并对Q2、Q3统一探索受控列生成与更强资源割。

\subsection{推广价值}
本文框架可推广至海岛医疗运输、异构无人机配送和带中途补能的应急物资运输。只需替换车辆参数、补能规则、任务优先级和时间窗，即可复用路线仿真、候选生成与排班模块。

\section{结论}
本文围绕海上油田人员运输建立了逐层递进的统一优化框架。Q1采用容量动态规划、广义节约、关联度引导ALNS、精英路线池主问题与标准ALNS教育，以89个架次覆盖1600人，总飞机使用时间14730 min，相对广义节约下降6.435\%，相对原问题理论下界的保守区间为10.445\%。Q2采用候选路线RMP、4路线精确局部MILP、长列增广、质量精英重组、全局重启、检查点续算与低频5路线强化，以95个架次覆盖4000人，总飞机使用时间17218 min，相对分开运输基线下降61.031\%，燃油下降62.577\%，座位利用率达到92.703\%；历史RMP界和局部邻域间隙均严格限定证明范围。Q3复用燃油可行路线并采用P0两阶段反馈、12启动Top-$3$、硬任务骨架--柔性Regret构造、时间感知$k\to m$重构、Closure--P2重构、PRO V1精英池强化、PRO V2异构筛选与多岛深度ALNS、Optional Rescue与强制P0反馈，以162个架次完成3840名非临时人员，将第一阶段基准从30546降至28728 min；分层连续多商品流给出全局下界14125 min，对应认证区间50.832\%。第二阶段在同一上界下服务157名临时人员，尚差3人达到理论上界160，157为固定架次结构下的可证明最优。三问最终提交方案均经独立Validator验证为零违规。

% ================================================================
% 官方要求：AI工具使用声明必须在参考文献之前，二选一并如实填写。
% ================================================================
\section*{AI 工具使用声明}
本参赛队在竞赛过程中使用了 AI 工具，主要用于算法思路梳理、LaTeX排版与语言润色、
Python代码调试、测试用例构造和图表生成；模型选择、程序运行、结果核验与论文定稿均由
参赛队员人工完成，详细使用情况见支撑材料。

% 若完全未使用AI，请删除上一句并改为：
% 本参赛队在竞赛过程中未使用任何 AI 工具。

\addcontentsline{toc}{section}{参考文献}
\bibliographystyle{unsrt}
\bibliography{references}

% 用于人工检查“正文+AI声明+参考文献”是否超过30页。
\label{body-end}
\clearpage
\appendix

\section{支撑材料文件列表}
\begin{longtable}{p{0.38\linewidth}p{0.52\linewidth}}
  \caption{支撑材料文件清单}\label{tab:support-files}\\
  \toprule
  文件或目录 & 内容与用途 \\
  \midrule
  \endfirsthead
  \toprule 文件或目录 & 内容与用途 \\ \midrule
  \endhead
  \path{code/ALNS_Q1_README.md} & 问题一环境、依赖、运行命令、两阶段参数和输出路径 \\
  \texttt{AI工具使用详情.md} & 工具名称/型号、使用环节、提示方式、采纳修改及人工核验情况 \\
  \path{code/requirements.txt} & Python 依赖与版本 \\
  \path{code/scripts/04_solve_q1_alns.py} & 问题一统一运行入口 \\
  \path{code/scripts/05_run_alns_multiseed.py} & 问题一关联度与标准ALNS公平多种子试验 \\
  \path{code/scripts/06_finalize_q1_relatedness.py} & 问题一融合最终解固化入口 \\
  \path{code/src/solver/q1_or.py} & 问题一精英路线池、精确主问题、LP放松与定向修复 \\
  \path{code/scripts/07_run_q1_route_pool_master.py}--\path{code/scripts/13_finalize_q1_or.py} & 问题一路线池强化、对照审计与最终固化入口 \\
  \path{code/scripts/05_solve_q2.py} & 问题二统一运行入口 \\
  \path{code/scripts/06_optimize_q2_lns.py} & 问题二精确局部MILP与自适应LNS入口 \\
  \path{code/scripts/07_recombine_q2_elites.py} & 问题二质量约束精英重组入口 \\
  \path{code/src/solver/q2_round3.py} & 问题二第三轮吸收审计、盆地指纹与局部主问题队列 \\
  \path{code/scripts/11_audit_q2_round3.py}--\path{code/scripts/13_summarize_q2_round3.py} & 问题二第三轮审计、吸收搜索与汇总入口 \\
  \path{code/scripts/06_solve_q3.py} & 问题三完整两阶段深度求解入口 \\
  \path{code/scripts/07_finalize_q3_results.py} & 问题三结果固化、两阶段校验和界限报告入口 \\
  \path{code/scripts/08_compute_q3_enhanced_bounds.py} & 问题三增强全局下界与候选池参考界计算入口 \\
  \path{code/scripts/09_optimize_q3_p0_p1.py} & 问题三投影反馈、多启动和多航班重构统一入口 \\
  \path{code/scripts/10_resume_q3_p2.py} & 问题三从Closure检查点继续执行P2的入口 \\
  \path{code/scripts/11_finalize_q3_p2_feedback.py} & 问题三最终第二阶段投影反馈与结果固化入口 \\
  \path{code/src/solver/q3_bounds.py} & 机场--机型分层连续多商品流和候选池连续主问题 \\
  \path{code/src/solver/q3_timing.py} & 允许海上等待的时间感知整数分钟排班器 \\
  \path{code/src/solver/q3_closure_p2.py} & Closure Gate、广义$k\to m$、P2A--P2D与缓存审计 \\
  \path{code/src/solver/} & Q1--Q3候选生成、ALNS/MILP、排班、破坏--修复和路线评价程序 \\
  \path{code/src/validation/} & 独立Validator \\
  \path{code/tests/} & 边界测试、属性测试与回归测试 \\
  \path{code/configs/} & 题目参数与机型参数 \\
  \path{code/outputs/q1/final/} & Q1融合最终CSV、配置、日志与零违规报告 \\
  \path{code/outputs/q2/best/} & Q2最终CSV、指标、图表数据与零违规报告 \\
  \path{code/outputs/q2/separate/} & Q2三类分开运输基线与校验报告 \\
  \path{code/outputs/q3/best/} & Q3 PRO V1阶段最佳解（29155 min）产物 \\
  \path{code/outputs/q3/closure_p2_best/bounds.json} & Closure--P2阶段增强界限 \\
  \path{code/outputs/q3/closure_p2_best/} & Closure--P2审计、消融轨迹、最终CSV和两阶段零违规报告 \\
  \path{code/outputs/q3/q3-pro/} & 问题三PRO V1精英池邻域搜索运行产物 \\
  \path{code/outputs/q3/q3-pro-v2/current_incumbent/} & Q3正式提交CSV、指标、认证界限（下界14125/上界28728/区间50.832\%）与两阶段零违规报告 \\
  \path{code/outputs/q3/q3-pro-v2/reports/} & PRO V2参数筛选、四岛深度ALNS、Optional Rescue与强制P0反馈轨迹报告 \\
  \path{code/src/solver/q3_pro.py} & 问题三PRO V1精英池路线库定价与结构强化 \\
  \path{code/src/solver/q3_pro_v2.py} & 问题三PRO V2异构筛选、多岛深度ALNS与救援邻域 \\
  \path{code/scripts/12_run_q3_pro.py}--\path{code/scripts/14_run_q3_pro_v2_robustness.py} & 问题三PRO V1/V2运行与稳健性检验入口 \\
  \texttt{figures/} & 论文图表的可复现CSV、PDF与PNG文件 \\
  \bottomrule
\end{longtable}

\section{完整源程序代码}
按照官方要求，附录应包含建模所用全部完整、可运行的源程序。下列路径应与支撑材料保持一致；若代码篇幅过大，可通过减小代码字号控制附录体积，但不要以伪代码替代完整源程序。

\inputcode{问题一两阶段ALNS运行入口}{code/scripts/04_solve_q1_alns.py}
\inputcode{问题一单设施容量动态规划}{code/src/solver/baseline.py}
\inputcode{问题一广义节约初解}{code/src/solver/improve.py}
\inputcode{问题一ALNS与整数规划修复}{code/src/solver/alns.py}
\inputcode{问题一关联度与Context候选排序}{code/src/solver/relatedness.py}
\inputcode{问题一多种子公平试验入口}{code/scripts/05_run_alns_multiseed.py}
\inputcode{问题一融合结果固化入口}{code/scripts/06_finalize_q1_relatedness.py}
\inputcode{问题一精英路线池与精确主问题}{code/src/solver/q1_or.py}
\inputcode{问题一路线池主问题入口}{code/scripts/07_run_q1_route_pool_master.py}
\inputcode{问题一定向强化入口}{code/scripts/08_run_q1_targeted_intensification.py}
\inputcode{问题一精英重组入口}{code/scripts/09_run_q1_elite_recombination.py}
\inputcode{问题一路径重连入口}{code/scripts/10_run_q1_path_relink.py}
\inputcode{问题一修复算法对照入口}{code/scripts/11_run_q1_repair_challengers.py}
\inputcode{问题一学习数据审计入口}{code/scripts/12_build_q1_ml_foundation.py}
\inputcode{问题一最终OR结果固化入口}{code/scripts/13_finalize_q1_or.py}
\inputcode{技术经停与燃油标签搜索}{code/src/solver/technical_stops.py}
\inputcode{共享求解缓存}{code/src/solver/cache.py}
\inputcode{预计算航段物理量}{code/src/solver/physics.py}
\inputcode{路线精确评价器}{code/src/solver/evaluator.py}
\inputcode{求解数据结构}{code/src/solver/models.py}
\inputcode{结果文件导出}{code/src/solver/exporter.py}
\inputcode{独立可行性校验器}{code/src/validation/solution_validator.py}
\inputcode{中文收敛图生成}{code/scripts/plot_q1_alns.py}
\inputcode{问题二候选路线整数主问题}{code/src/solver/q2.py}
\inputcode{问题二统一运行入口}{code/scripts/05_solve_q2.py}
\inputcode{问题二精确局部MILP与自适应LNS}{code/src/solver/q2_lns.py}
\inputcode{问题二有向OD流候选生成}{code/src/solver/q2_flow.py}
\inputcode{问题二精英池与差异重组}{code/src/solver/q2_round2.py}
\inputcode{问题二第三轮吸收式强化}{code/src/solver/q2_round3.py}
\inputcode{问题二局部搜索运行入口}{code/scripts/06_optimize_q2_lns.py}
\inputcode{问题二精英重组运行入口}{code/scripts/07_recombine_q2_elites.py}
\inputcode{问题二第三轮解审计}{code/scripts/11_audit_q2_round3.py}
\inputcode{问题二吸收式邻域搜索}{code/scripts/12_search_q2_absorption.py}
\inputcode{问题二第三轮结果汇总}{code/scripts/13_summarize_q2_round3.py}
\inputcode{问题二结果图生成}{code/scripts/plot_q2_results.py}
\inputcode{问题一二融合阶段图生成}{code/scripts/plot_integrated_improvements.py}
\inputcode{问题三时间窗排班与深度邻域}{code/src/solver/q3.py}
\inputcode{问题三完整两阶段求解入口}{code/scripts/06_solve_q3.py}
\inputcode{问题三结果固化与界限校验}{code/scripts/07_finalize_q3_results.py}
\inputcode{问题三增强求解界限模型}{code/src/solver/q3_bounds.py}
\inputcode{问题三增强求解界限入口}{code/scripts/08_compute_q3_enhanced_bounds.py}
\inputcode{问题三时间感知排班器}{code/src/solver/q3_timing.py}
\inputcode{问题三P0--P1反馈与结构优化入口}{code/scripts/09_optimize_q3_p0_p1.py}
\inputcode{问题三Closure--P2核心算法}{code/src/solver/q3_closure_p2.py}
\inputcode{问题三P2检查点续算入口}{code/scripts/10_resume_q3_p2.py}
\inputcode{问题三最终反馈与结果固化}{code/scripts/11_finalize_q3_p2_feedback.py}
\inputcode{问题三PRO V1精英池强化}{code/src/solver/q3_pro.py}
\inputcode{问题三PRO V1运行入口}{code/scripts/12_run_q3_pro.py}
\inputcode{问题三PRO V2深度强化}{code/src/solver/q3_pro_v2.py}
\inputcode{问题三PRO V2运行入口}{code/scripts/13_run_q3_pro_v2.py}
\inputcode{问题三PRO V2稳健性检验入口}{code/scripts/14_run_q3_pro_v2_robustness.py}
\inputcode{问题三中文图表生成}{code/scripts/plot_q3_results.py}
\inputcode{问题三算法演进图生成}{code/scripts/plot_q3_search_evolution.py}

\section{补充结果与校验报告}
问题一完整结果位于\path{code/outputs/q1/final/}；问题二完整路线、4000条人员分配、
方法对比、代表性路线状态和独立校验报告位于\path{code/outputs/q2/best/}，分开运输基线
位于\path{code/outputs/q2/separate/}。问题三正式提交文件位于\path{code/outputs/q3/q3-pro-v2/current_incumbent/}，
筛选、四岛深度ALNS、Rescue与反馈轨迹报告位于\path{code/outputs/q3/q3-pro-v2/reports/}，
Closure--P2阶段CSV、界限、投影审计、消融轨迹和校验报告位于
\path{code/outputs/q3/closure_p2_best/}，PRO V1阶段产物位于\path{code/outputs/q3/best/}。

\section{提交前合规核对}
\begin{enumerate}[label=$\square$]
  \item 摘要、标题和关键词严格位于第 1 页；正文从第 2 页开始；全文无目录。
  \item 正文、AI 声明和参考文献合计不超过 30 页；附录页数不限。
  \item 全文无姓名、学校、指导教师、账号和本机用户名等身份信息。
  \item Q1/Q2/Q3 的人员行数、路线编号、站序、容量、油量和时间均通过最终校验。
  \item Q3 最终总飞机使用时间不超过同一口径计算的 $T_0$。
  \item 摘要、结果表、正文和最终 CSV 中的五项指标完全一致。
  \item AI 声明如实填写，支撑材料含\texttt{AI工具使用详情.md}。
  \item 论文 PDF 和支撑材料压缩包分别按三位参赛队号命名，且均不超过 20 MB。
  \item 全文无草稿标记，并在干净目录中完成一次一键复现。
\end{enumerate}

\end{document}
